%%%%%%%%%%%%%%%%%%%%%%%%%%%%%%%%%%%%%%%%%%%%%%%%%%%%%%%%%%%%%%%%%%%%%%%%%%%%%
%  LINE OPERATORS — EXTENSIONS
%  Module categories, dg-shifted Yangian, MC3 programme, W-algebra lines
%%%%%%%%%%%%%%%%%%%%%%%%%%%%%%%%%%%%%%%%%%%%%%%%%%%%%%%%%%%%%%%%%%%%%%%%%%%%%

%%%%%%%%%%%%%%%%%%%%%%%%%%%%%%%%%%%%%%%%%%%%%%%%%%%%%%%%%%%%%%%%%%%%%%%%%%%%%
% SECTION — MODULE CATEGORIES OVER THE KOSZUL DUAL
%%%%%%%%%%%%%%%%%%%%%%%%%%%%%%%%%%%%%%%%%%%%%%%%%%%%%%%%%%%%%%%%%%%%%%%%%%%%%
\section{Module categories over the Koszul dual}
% label removed: sec:module-categories-koszul-dual
\index{Koszul duality!module categories|textbf}

Theorem~\ref{thm:lines_as_modules} identifies
$\cC_{\mathrm{line}}\simeq\cA^!\text{-}\mathbf{mod}$
at the level of objects and morphisms.  This section
promotes the identification to the full dg-categorical
level: enriched hom complexes, $A_\infty$-module structures,
and the monoidal product induced by the $R$-matrix.
Throughout, $\cA$ denotes a chiral Koszul algebra satisfying
(H1)--(H4), with bar complex
$\barB_{\mathrm{ch}}(\cA)$ and Koszul dual
$\cA^!=H^*(\barB_{\mathrm{ch}}(\cA))^\vee$
(Construction~\ref{constr:dg-shifted-yangian-from-bar}).

\begin{remark}[Module structures: $A_\infty$ first, strict second]
% label removed: rem:module-ainfty-first
\index{module category!A-infinity@$A_\infty$ structure}
The $\cA^!$-module structure on a line-operator state space
$V_\ell$ is an $A_\infty$-module in the sense of
Keller~\cite{Kel06}: a sequence of multilinear maps
$\mu_n^M\colon(\cA^!)^{\otimes(n-1)}\otimes V_\ell\to V_\ell$
($n\ge 1$) satisfying the $A_\infty$-module relations
$\sum_{i+j=n+1}\mu_i^M\circ(\mathrm{id}^{\otimes(i-j)}\otimes\mu_j^M)=0$.
The strict dg module obtained by retaining only $\mu_1^M$ and
$\mu_2^M$ is a model in the sense of
Convention~\ref{conv:vol2-strict-models}: it computes the
correct derived category but not the correct $A_\infty$-enrichment.
\end{remark}

\subsection{The dg category $\cA^!\text{-}\mathbf{mod}$}
% label removed: subsec:dg-category-amod

\begin{definition}[$A_\infty$-modules over $\cA^!$]
% label removed: def:ainfty-module-koszul-dual
\index{A-infinity module@$A_\infty$-module!over Koszul dual}
An \emph{$A_\infty$-module} over the dg algebra $\cA^!$
is a graded vector space $M$ equipped with
structure maps $\mu_n^M\colon(\cA^!)^{\otimes(n-1)}\otimes M\to M$
of degree $2-n$ for all $n\ge 1$, satisfying:
\begin{equation}
% label removed: eq:ainfty-module-relations
\sum_{\substack{i+j=n+1 \\ i\ge 1,\; j\ge 1}}
(-1)^{i(j-1)+\sum_{k=1}^{i-1}|a_k|}
\mu_i^M(a_1,\ldots,a_{i-1},
\mu_j^M(a_i,\ldots,a_{i+j-2},m))
\;=\;0,
\end{equation}
where $|a_k|$ denotes the cohomological degree of $a_k\in\cA^!$
and $m\in M$.  Setting $n=1$: $(\mu_1^M)^2=0$, so $\mu_1^M$
is the internal differential of~$M$.  Setting $n=2$:
$\mu_1^M\circ\mu_2^M+\mu_2^M\circ(\mu_1^{\cA^!}\otimes\mathrm{id}+\mathrm{id}\otimes\mu_1^M)=0$,
which is the Leibniz rule for the action $\mu_2^M$.
Higher $\mu_n^M$ ($n\ge 3$) encode homotopy-coherent corrections:
the action of $\cA^!$ on $M$ is associative only up to
specified higher homotopies.
\end{definition}

\begin{definition}[Enriched hom complex]
% label removed: def:enriched-hom
\index{enriched hom complex|textbf}
For $A_\infty$-modules $M$ and $N$ over $\cA^!$,
the \emph{enriched hom complex}
$\underline{\Hom}_{\cA^!}(M,N)$ is the chain complex
with
\begin{equation}
% label removed: eq:enriched-hom
\underline{\Hom}_{\cA^!}^k(M,N)
\;:=\;
\prod_{n\ge 0}
\Hom_\Bbbk\bigl((\cA^!)^{\otimes n}\otimes M,\;N\bigr)^k,
\end{equation}
where the degree-$k$ component consists of families of
linear maps $f_n\colon(\cA^!)^{\otimes n}\otimes M\to N$
of degree $k-n$ for each $n\ge 0$.  The differential
$\partial$ on $\underline{\Hom}_{\cA^!}(M,N)$ is defined by
\begin{align}
% label removed: eq:enriched-hom-diff
(\partial f)_n(a_1,\ldots,a_n,m)
\;=\;&\;
\sum_{j=1}^{n}(-1)^{\epsilon_j}
f_n(a_1,\ldots,\mu_1^{\cA^!}(a_j),\ldots,a_n,m) \notag \\
&+\;(-1)^{\epsilon_{n+1}}f_n(a_1,\ldots,a_n,\mu_1^M(m)) \notag \\
&-\;(-1)^{|f|}\mu_1^N(f_n(a_1,\ldots,a_n,m)) \notag \\
&+\;\sum_{\substack{i+j=n+1 \\ i\ge 2}}
(-1)^{\epsilon_i}
f_{n-j+1}(a_1,\ldots,a_{i-1},
\mu_j^M(a_i,\ldots,a_{i+j-2},m)) \notag \\
&-\;\sum_{\substack{i+j=n+1 \\ j\ge 2}}
(-1)^{|f|+\epsilon_i}
\mu_j^N(a_1,\ldots,a_{j-1},
f_{n-j+1}(a_j,\ldots,a_n,m)),
\end{align}
where $\epsilon_j=\sum_{k=1}^{j-1}|a_k|$.  Then
$\partial^2=0$ by the $A_\infty$-module relations
\eqref{eq:ainfty-module-relations} for both $M$ and $N$.
\end{definition}

\begin{proposition}[$\partial^2=0$ on the enriched hom;
\ClaimStatusProvedHere]
% label removed: prop:enriched-hom-d-squared
\index{enriched hom complex!differential}
The differential $\partial$ of
Definition~\textup{\ref{def:enriched-hom}} satisfies
$\partial^2=0$.
\end{proposition}

\begin{proof}
Write $\partial=\partial_{\mathrm{int}}+\partial_M+\partial_N$
corresponding to the three groups of terms in
\eqref{eq:enriched-hom-diff}: the internal differential
on $\cA^!$ and $M$, the $A_\infty$-module structure of $M$,
and the $A_\infty$-module structure of $N$.  Each of
$\partial_{\mathrm{int}}^2$, $\partial_M^2$, $\partial_N^2$
vanishes by the $A_\infty$-module relations.  The cross terms
$\partial_{\mathrm{int}}\partial_M+\partial_M\partial_{\mathrm{int}}$
vanish by the Leibniz component ($n=2$) of the
$A_\infty$-module relation for $M$, and similarly for the
$N$ cross terms.  The mixed term
$\partial_M\partial_N+\partial_N\partial_M$ vanishes because
$f$ intertwines the $M$- and $N$-module structures:
explicitly, the sum over partitions $i+j+k=n+2$
telescopes to zero by the signed Koszul convention on
triple compositions.  The full $\partial^2$ therefore
decomposes into six pairs, each vanishing by the
$A_\infty$-module identities.
\end{proof}

\begin{remark}[Topology on the enriched hom complex]
% label removed: rem:enriched-hom-topology
\index{enriched hom complex!topology}
The enriched hom complex~\eqref{eq:enriched-hom} carries
the product topology: the projective limit over $n$ of
the finite-dimensional spaces
$\Hom_\Bbbk((\cA^!)^{\otimes n}\otimes M,\,N)^k$.  The
differential $\partial$
of~\eqref{eq:enriched-hom-diff} converges in this topology
because each arity-$n$ component $(\partial f)_n$ depends
only on $f_{n-1}$, $f_n$, and $f_{n+1}$: the sum
$\sum_{i+j=n+1}$ has finitely many terms at each fixed
arity.  Consequently the series defining $\partial$
stabilizes arity-by-arity, and no analytic convergence
condition is needed.  For the same reason, the composition
$(g\circ f)_n = \sum_{i+j=n}g_i\circ(\mathrm{id}^{\otimes i}
\otimes f_j)$ is a finite sum at each arity and converges
in the product topology.
\end{remark}

\begin{definition}[The dg category $\cA^!\text{-}\mathbf{mod}$]
% label removed: def:dg-category-amod
\index{module category!dg category|textbf}
The \emph{dg category} $\cA^!\text{-}\mathbf{mod}$
has:
\begin{itemize}
\item \emph{Objects}: $A_\infty$-modules $(M,\{\mu_n^M\}_{n\ge 1})$
  over $\cA^!$ (Definition~\ref{def:ainfty-module-koszul-dual});
\item \emph{Morphism complexes}:
  $\Hom_{\cA^!\text{-}\mathbf{mod}}(M,N)
  :=\underline{\Hom}_{\cA^!}(M,N)$
  (Definition~\ref{def:enriched-hom});
\item \emph{Composition}: for $f\in\underline{\Hom}(M,N)$
  and $g\in\underline{\Hom}(N,P)$, the composite
  $(g\circ f)_n:=
  \sum_{i+j=n}g_i\circ(\mathrm{id}^{\otimes i}\otimes f_j)$;
\item \emph{Identities}: the family $(\mathrm{id}_M)_0=\mathrm{id}_M$,
  $(\mathrm{id}_M)_n=0$ for $n\ge 1$.
\end{itemize}
The composition is associative and compatible with the
differential $\partial$ by the standard bar-cobar argument
(cf.~\cite{LH03}, \S3.2).
\end{definition}

\begin{theorem}[dg-categorical equivalence;
\ClaimStatusProvedHere]
% label removed: thm:dg-categorical-equivalence
\index{line operators!dg-categorical equivalence|textbf}
Let\/ $\cA$ be a logarithmic\/ $\SCchtop$-algebra.  The functor
$\Phi\colon\cC_{\mathrm{line}}\to\cA^!\text{-}\mathbf{mod}$
of Theorem~\textup{\ref{thm:lines_as_modules}} is a
quasi-equivalence of dg categories: for all line operators
$\ell_1,\ell_2$, the induced map
\begin{equation}
% label removed: eq:dg-cat-qiso
\Phi_{\ell_1,\ell_2}\colon
\Hom_{\cC_{\mathrm{line}}}(\ell_1,\ell_2)
\;\xrightarrow{\;\sim\;}
\underline{\Hom}_{\cA^!}(V_{\ell_1},V_{\ell_2})
\end{equation}
is a quasi-isomorphism of chain complexes.
\end{theorem}

\begin{proof}
\textbf{Step 1: The bar--cobar adjunction at the module level.}
The Quillen equivalence
$\Omega\colon\mathrm{Coalg}_{\mathbf{B}}
\rightleftarrows
\mathrm{Alg}_{\SCchtop}\colon\mathbf{B}$
(Theorem~\ref{thm:bar-cobar-adjunction}) extends to module
categories by the standard operadic transfer
(\cite{BM03}, Theorem 4.7): there is a Quillen
equivalence
\begin{equation}
% label removed: eq:module-quillen-equiv
\Omega^{\mathrm{mod}}\colon
\mathrm{Comod}_{\mathbf{B}(\cA)}
\rightleftarrows
\cA\text{-}\mathrm{Mod}_{\SCchtop}
\colon\mathbf{B}^{\mathrm{mod}},
\end{equation}
where $\cA\text{-}\mathrm{Mod}_{\SCchtop}$ denotes the
category of $\SCchtop$-modules over the $\SCchtop$-algebra
$\cA$, and $\mathrm{Comod}_{\mathbf{B}(\cA)}$ denotes
conilpotent comodules over the bar coalgebra.

The homotopy-Koszulity of $\SCchtop$
(Theorem~\ref{thm:homotopy-Koszul}) guarantees this is a
Quillen equivalence (not merely an adjunction).  The
argument is as follows.  Berger--Moerdijk \cite{BM03},
Theorem~4.7 requires the operad to be Koszul for the
module-level bar-cobar adjunction to be a Quillen
equivalence.  For homotopy-Koszul operads, the same
conclusion holds by the two-out-of-three property: if
$\varphi\colon\cO\xrightarrow{\sim}\cO'$ is a
quasi-isomorphism of operads and $\cO'$ is Koszul,
then the induced functor on module categories is a
Quillen equivalence (because the bar-cobar adjunction
for $\cO$ factors through the one for $\cO'$ via the
quasi-isomorphism $\varphi$, and both the bar and cobar
functors preserve quasi-isomorphisms between cofibrant
objects).  Homotopy-Koszulity of $\SCchtop$
(Theorem~\ref{thm:homotopy-Koszul}) provides precisely
such a quasi-isomorphism
$\SCchtop\xrightarrow{\sim}\mathrm{SC}$ to the
classical Swiss-cheese operad, and $\mathrm{SC}$ is
Koszul (Livernet~\cite{Liv15}; the binary quadratic
case follows from Ginzburg--Kapranov~\cite{GK94}).

\textbf{Step 2: From comodules to dual modules.}
For a conilpotent coalgebra $C$ over a field, the
linear duality functor $(-)^\vee\colon C\text{-Comod}
\xrightarrow{\sim}C^\vee\text{-Mod}$ is an equivalence of
dg categories (this is standard: conilpotence ensures every
comodule is the union of its finite-dimensional subcomodules,
and dualizing commutes with filtered colimits over a field).
Applied to $C=\barB_{\mathrm{ch}}(\cA)$:
\[
\mathrm{Comod}_{\barB_{\mathrm{ch}}(\cA)}
\;\simeq\;
\barB_{\mathrm{ch}}(\cA)^\vee\text{-Mod}.
\]
Passing to cohomology (which is an equivalence of dg
categories because $\barB_{\mathrm{ch}}(\cA)^\vee$ is
formal by chiral Koszulness of $\cA$), we obtain
$\barB_{\mathrm{ch}}(\cA)^\vee\text{-Mod}
\simeq\cA^!\text{-}\mathbf{mod}$.

\textbf{Step 3: dg-categorical enrichment.}
Combining Steps 1 and 2, we have a chain of Quillen
equivalences
\[
\cA\text{-}\mathrm{Mod}_{\SCchtop}
\;\xleftarrow{\;\sim\;}
\mathrm{Comod}_{\mathbf{B}(\cA)}
\;\xrightarrow{\;\sim\;}
\cA^!\text{-}\mathbf{mod}.
\]
The physical line-operator category $\cC_{\mathrm{line}}$
is, by construction, the full dg subcategory of
$\cA\text{-}\mathrm{Mod}_{\SCchtop}$ on the physically
realizable modules (those arising from boundary conditions
under (H1)--(H4)).  By the argument of
Theorem~\ref{thm:lines_as_modules}, Step 4 (essential
surjectivity), every $\cA^!$-module is realized; hence
$\cC_{\mathrm{line}}$ is quasi-equivalent to the entire
module category $\cA^!\text{-}\mathbf{mod}$.

The map on morphism complexes is the composite
\begin{align*}
\Hom_{\cC_{\mathrm{line}}}(\ell_1,\ell_2)
&\;\xrightarrow{\;\mathbf{B}^{\mathrm{mod}}\;}
\Hom_{\mathrm{Comod}}(\mathbf{B}^{\mathrm{mod}}V_{\ell_1},
\mathbf{B}^{\mathrm{mod}}V_{\ell_2}) \\
&\;\xrightarrow{\;(-)^\vee\;}
\Hom_{\cA^!\text{-mod}}(V_{\ell_1}^\vee,V_{\ell_2}^\vee)
\;\cong\;
\underline{\Hom}_{\cA^!}(V_{\ell_1},V_{\ell_2}).
\end{align*}
Each arrow is a quasi-isomorphism (the first by the
Quillen equivalence, the second by conilpotent duality),
so $\Phi_{\ell_1,\ell_2}$ is a quasi-isomorphism.
\end{proof}

\begin{remark}[Strict versus $A_\infty$ modules]
% label removed: rem:strict-vs-ainfty-modules
\index{strict model!module categories}
The derived category $D(\cA^!)$ (= the homotopy category
of strict dg modules) captures the same triangulated
category as $H^0(\cA^!\text{-}\mathbf{mod})$.  However,
the full $A_\infty$-enrichment records higher information:
Massey products in $\operatorname{Ext}^*_{\cA^!}(M,N)$,
obstructions to formality of the module category, and
the $A_\infty$-Yoneda embedding
$\cA^!\text{-}\mathbf{mod}\hookrightarrow
\mathrm{Fun}(\cA^!,\mathrm{Ch})$.  The strict model is
recovered by the minimal model theorem: every
$A_\infty$-module is $A_\infty$-quasi-isomorphic to a
minimal one (where $\mu_1^M=0$), and the resulting
minimal $A_\infty$-module category is the Kadeishvili
$A_\infty$-structure on $\operatorname{Ext}^*_{\cA^!}$
(\cite{Kad82}, Convention~\ref{conv:vol2-strict-models}).
\end{remark}

\subsection{Explicit module categories for standard families}
% label removed: subsec:explicit-module-categories

We extract the $A_\infty$-module structures for the three
principal families.

\subsubsection{Heisenberg: curved symmetric modules}
% label removed: subsubsec:heisenberg-modules

For $\cA=\cH_k$ (Heisenberg at level $k$),
the Koszul dual is $\cA^!=\mathrm{Sym}^{\mathrm{ch}}(V^*)$:
the chiral symmetric algebra on the dual (Vol~I, not $\cH_{-k}$).  The bar complex
$\barB_{\mathrm{ch}}(\cH_k)$ is the symmetric coalgebra
$\mathrm{Sym}^{\mathrm{ch}}(V[-1])$ with $V=\Bbbk\cdot J$
the one-dimensional current space, and the bar differential
reduces to the coproduct of $\mathrm{Sym}$.  The Koszul
dual is
\[
\cH_{-k}
\;=\;
\bigl(\mathrm{Sym}^{\mathrm{ch}}(V[-1])\bigr)^\vee
\;\cong\;
\mathrm{Sym}^{\mathrm{ch}}(V^*[1]),
\]
with the dual inner product $\langle J^*,J^*\rangle=-k$.

\begin{proposition}[Heisenberg modules as twisted Fock spaces;
\ClaimStatusProvedHere]
% label removed: prop:heisenberg-modules-fock
\index{Heisenberg algebra!modules!twisted Fock}
The $A_\infty$-module category
$\cH_{-k}\text{-}\mathbf{mod}$ is equivalent (as a
dg category) to the category of $\Z$-graded vector spaces:
\begin{equation}
% label removed: eq:heisenberg-module-category
\cH_{-k}\text{-}\mathbf{mod}
\;\simeq\;
\mathrm{Vect}^\Z.
\end{equation}
The simple modules are the Fock spaces
$\cF_q$ ($q\in\Z$), with $\cH_{-k}$-action:
\begin{itemize}
\item $J^*_{(0)}\cdot v=q\cdot v$ for $v\in\cF_q$
  (charge eigenvalue);
\item $J^*_{(n)}\cdot v=0$ for $n>0$ (vacuum condition);
\item $J^*_{(-n)}$ ($n>0$) acts freely as creation operators.
\end{itemize}
The higher $A_\infty$-module maps $\mu_n^M$ vanish
for $n\ge 3$: the $\cH_{-k}$-module structure is
strictly associative, reflecting the Gaussianity
(shadow depth $r_{\max}=2$) of the Heisenberg algebra
(Vol~I, \S\ref*{sec:heisenberg-shadow-gaussianity}).
\end{proposition}

\begin{proof}
The Heisenberg algebra $\cH_{-k}$ is generated by a single
current $J^*$ with OPE $J^*(z)J^*(w)\sim -k/(z-w)^2$.
Its module theory is that of a rank-$1$ lattice vertex algebra
with trivial lattice: every module is a direct sum of Fock
spaces $\cF_q$ indexed by the charge $q$ (the eigenvalue of
$J^*_{(0)}$).  The category of such modules is
$\mathrm{Vect}^\Z$, with morphisms between distinct charges
being zero and endomorphisms of $\cF_q$ being the identity.

The vanishing of $\mu_n^M$ for $n\ge 3$ follows from the
homological perturbation lemma (HPL).  By the HPL, the
transferred $A_\infty$-module structure has
$\mu_n^M = \sum_T \mu_T^M$, where the sum is over planar
trees $T$ with $n$ inputs, and each internal vertex of $T$
is labelled by a bar cohomology class.  For the Heisenberg,
bar cohomology is concentrated in degree~$1$
(shadow depth $r_{\max}=2$; Vol~I,
\S\ref*{sec:heisenberg-shadow-gaussianity}): the only
nonzero bar cohomology class is the degree-$1$ class
$[s^{-1}J]\in H^1(\barB(\cH_k))$.  The transferred
operations $\mu_n^M$ for $n\ge 3$ require at least one
internal vertex labelled by a bar cohomology class of
degree $\ge 2$ (to produce a tree with $n$ inputs from
binary vertices and homotopy propagators), but no such
classes exist.  Therefore the only nonzero transferred
operations are $\mu_1^M$ (the differential) and $\mu_2^M$
(the binary action), and $\mu_n^M = 0$ for all $n\ge 3$.
Equivalently, $\cH_{-k}$ is intrinsically formal as an
$A_\infty$-algebra (it is a quadratic algebra with
quadratic relations), so its $A_\infty$-module structures
are determined by the binary action $\mu_2^M$.
\end{proof}

\begin{computation}[Enriched hom for Heisenberg Fock modules;
\ClaimStatusProvedHere]
% label removed: comp:enriched-hom-heisenberg
\index{enriched hom complex!Heisenberg}
For $q_1\ne q_2$:
$\underline{\Hom}_{\cH_{-k}}(\cF_{q_1},\cF_{q_2})=0$
(the charge grading is strict).  For $q_1=q_2=q$:
\begin{equation}
% label removed: eq:hom-heisenberg
\underline{\Hom}_{\cH_{-k}}(\cF_q,\cF_q)
\;\cong\;
\Bbbk
\end{equation}
concentrated in degree $0$ (the identity intertwiner).
All Ext groups vanish:
$\operatorname{Ext}^n_{\cH_{-k}}(\cF_{q_1},\cF_{q_2})=0$
for $n\ge 1$ and all $q_1,q_2$.  The dg category
$\cH_{-k}\text{-}\mathbf{mod}$ is formal.
\end{computation}

\subsubsection{Affine $\mathfrak{sl}_2$: Weyl modules and quantum group}
% label removed: subsubsec:affine-sl2-modules

For $\cA=\widehat{\mathfrak{sl}}_2{}_k$:
$\cA^!=\widehat{\mathfrak{sl}}_2{}_{-k-4}$
(level duality: $k\mapsto -k-2h^\vee=-k-4$ for
$\mathfrak{sl}_2$).

\begin{definition}[Weyl modules for the dual affine algebra]
% label removed: def:weyl-modules-dual
\index{Weyl module!dual affine algebra}
For $\lambda\in\Z_{\ge 0}$, the \emph{Weyl module}
$W_\lambda$ for $\widehat{\mathfrak{sl}}_2{}_{-k-4}$
is the universal highest-weight module generated by a
vector $v_\lambda$ satisfying:
\begin{enumerate}[label=(\roman*)]
\item $h_{(0)}\cdot v_\lambda=\lambda\cdot v_\lambda$
  (weight $\lambda$);
\item $e_{(n)}\cdot v_\lambda=0$ for $n\ge 0$
  (highest weight);
\item $f_{(n)}\cdot v_\lambda=0$ for $n\ge 1$
  (positive-mode annihilation);
\item $K\cdot v_\lambda=(-k-4)\cdot v_\lambda$
  (level $-k-4$).
\end{enumerate}
The Weyl module is the maximal integrable quotient of the
Verma module $M(\lambda,-k-4)$; when $\lambda\le k$
(integrable range for the original level~$k$), the
simple quotient $L_\lambda=W_\lambda/\mathrm{rad}(W_\lambda)$
is the irreducible highest-weight module.
\end{definition}

\begin{proposition}[Affine module category as quantum group
category; \ClaimStatusProvedHere]
% label removed: prop:affine-module-category
\index{Kac--Moody!module category!quantum group}
The dg category
$\widehat{\mathfrak{sl}}_2{}_{-k-4}\text{-}\mathbf{mod}$
carries a braided monoidal structure with braiding
$R(z)=\exp(\Omega/z(k{+}2))$, where
$\Omega=e\otimes f+f\otimes e+\tfrac12 h\otimes h$.
This braided monoidal category is equivalent to a
full subcategory of the shifted quantum group category:
\begin{equation}
% label removed: eq:affine-module-quantum-group
\widehat{\mathfrak{sl}}_2{}_{-k-4}\text{-}\mathbf{mod}
\;\simeq\;
\cO^{\mathrm{sh}}_q(\widehat{\mathfrak{sl}}_2),
\qquad
q=e^{i\pi/(k+2)}.
\end{equation}

The $A_\infty$-module maps $\mu_n^M$ ($n\ge 4$) vanish
on Weyl modules, reflecting shadow depth $r_{\max}=3$
(the affine algebra terminates at arity $3$;
Vol~I, \S\ref*{sec:affine-cubic-shadow}).  The non-vanishing
$\mu_3^M$ encodes the cubic shadow: on the tensor product
$W_{\lambda_1}\otimes W_{\lambda_2}\otimes W_{\lambda_3}$,
$\mu_3^M$ is the Massey product corresponding to the
non-trivial extension class
$\operatorname{Ext}^2_{\widehat{\fg}_{-k-4}}
(W_{\lambda_3},W_{\lambda_1}\otimes W_{\lambda_2})$.
\end{proposition}

\begin{proof}
The braided monoidal structure follows from the Yangian
coproduct $\Delta_z$ on $\cA^!=Y_\hbar(\mathfrak{sl}_2)$
(Theorem~\ref{thm:Koszul_dual_Yangian}).  The
$R$-matrix $R(z)=\exp(\Omega/z(k{+}2))$ is the
monodromy of the KZ connection
$\nabla=d-\Omega\,d\log(z)/(k+2)$
(Computation~\ref{comp:line-op-sl2}).  This monodromy
representation factors through the quantum group
$U_q(\mathfrak{sl}_2)$ at $q=e^{i\pi/(k+2)}$
(Drinfeld~\cite{Dri89}, Kazhdan--Lusztig~\cite{KL93}),
giving the equivalence with $\cO^{\mathrm{sh}}_q$.

For the $A_\infty$ claim: the shadow Postnikov tower
of $\widehat{\fg}_k$ terminates at arity $3$
(Vol~I, Theorem~\ref*{thm:affine-shadow-termination}).
In the module-theoretic language: the bar complex
$\barB_{\mathrm{ch}}(\widehat{\fg}_k)$ has
$H^n(\barB)\ne 0$ only for $n\le 3$, so the
$A_\infty$-structure on $\cA^!$ has $m_n=0$ for $n\ge 4$.
By transfer (the homological perturbation lemma), the
induced $A_\infty$-module maps $\mu_n^M$ on Weyl modules
also vanish for $n\ge 4$.

The non-vanishing $\mu_3^M$ is computed as follows.
Take three Weyl modules $W_{\lambda_1}$, $W_{\lambda_2}$,
$W_{\lambda_3}$ and the extension sequence
$0\to K\to W_{\lambda_1}\otimes W_{\lambda_2}
\to L_{\lambda_1+\lambda_2}\to 0$,
where $K$ is the kernel of the projection.  The Massey
product $\langle e_{(0)},f_{(0)},e_{(0)}\rangle$ in
$\operatorname{Ext}^*_{\widehat{\fg}_{-k-4}}$ is
non-trivial (it detects the Serre relation
$[e,[e,f]]=2e$ at the $A_\infty$ level), and this is
precisely $\mu_3^M$.
\end{proof}

\begin{computation}[Enriched hom for $\widehat{\mathfrak{sl}}_2$
Weyl modules; \ClaimStatusProvedHere]
% label removed: comp:enriched-hom-affine
\index{enriched hom complex!affine}
The enriched hom between Weyl modules computes the
$\fg$-equivariant cohomology of the configuration space
$\mathrm{Conf}_n(\C)$:
\begin{equation}
% label removed: eq:ext-affine
\operatorname{Ext}^n_{\widehat{\fg}_{-k-4}}
(W_\lambda,W_\mu)
\;\cong\;
\begin{cases}
\Bbbk & \text{if } n=0,\;\lambda=\mu, \\
\Hom_{\fg}(V_\lambda,V_\mu) & \text{if } n=0,\;
  \lambda\ne\mu, \\
H^n(\mathrm{Conf}_2(\C);
  \Hom_{\fg}(V_\lambda,V_\mu)_{\mathrm{loc}})
  & \text{if } n\ge 1,
\end{cases}
\end{equation}
where $V_\lambda$ denotes the finite-dimensional
$\mathfrak{sl}_2$-module of highest weight~$\lambda$,
and the local system
$\Hom_{\fg}(V_\lambda,V_\mu)_{\mathrm{loc}}$
is the monodromy local system of the KZ connection.
At generic level~$k$ (non-integral $q$):
$\operatorname{Ext}^n=0$ for $n\ge 1$, and the
module category is semisimple.  At integral $k$
(rational $q=e^{i\pi/(k+2)}$): the non-vanishing of
$\operatorname{Ext}^1$ between Weyl modules reflects
the non-semisimplicity of the Verlinde category.
\end{computation}

\subsubsection{Virasoro: highest-weight representations of $\mathrm{Vir}_{26-c}$}
% label removed: subsubsec:virasoro-modules

For $\cA=\mathrm{Vir}_c$:
$\cA^!=\mathrm{Vir}_{26-c}$
(Virasoro at central charge $26-c$).  This is the
depth-zero resonance shadow
(Vol~I, Remark~\ref*{rem:virasoro-resonance-model}): the
Koszul dual of the Virasoro algebra at central charge~$c$
is the Virasoro algebra at central charge $26-c$, not a
simple level shift.  The self-dual point is $c=13$
(not $c=26$).

\begin{proposition}[Virasoro module category;
\ClaimStatusProvedHere]
% label removed: prop:virasoro-module-category
\index{Virasoro algebra!module category}
The dg module category
$\mathrm{Vir}_{26-c}\text{-}\mathbf{mod}^{\mathrm{dg}}$
has:
\begin{itemize}
\item \emph{Simple objects}: Verma modules $M(h',26-c)$
  for $h'\in\C$, and their irreducible quotients $L(h',26-c)$
  (when singular vectors exist);
\item \emph{Morphism complexes}: the enriched hom
  $\underline{\Hom}_{\mathrm{Vir}_{26-c}}(M_1,M_2)$
  is computed by the BRST complex of the Virasoro algebra
  at central charge $26-c$;
\item \emph{$A_\infty$-structure}: non-trivial in all
  arities.
\end{itemize}
The $A_\infty$-module maps $\mu_n^M$ do \emph{not} vanish
for any finite~$n$: the Virasoro shadow tower is infinite
($r_{\max}=\infty$, class~M;
Vol~I, \S\ref*{sec:mixed-cubic-quartic-shadows}), so the
module category carries a genuinely infinite $A_\infty$
hierarchy.  This is the module-theoretic manifestation
of the mixed shadow depth of the Virasoro algebra.
\end{proposition}

\begin{proof}
The Verma module $M(h',26-c)$ is the standard Virasoro
highest-weight module with $L_0$-eigenvalue $h'$ and
central charge $26-c$.  Its irreducible quotient
$L(h',26-c)$ is obtained by quotienting by the maximal
proper submodule (generated by singular vectors at
levels determined by the Kac determinant formula at
central charge $26-c$).

For the $A_\infty$ claim: the algebra $\mathrm{Vir}_c$ has
$m_n\ne 0$ for all $n\ge 2$ (class~M, infinite shadow
depth; Vol~I, \S\ref*{sec:mixed-cubic-quartic-shadows}).
By the homological perturbation lemma, the transferred
$A_\infty$-module maps $\mu_n^M$ on a \emph{generic}
module are non-trivial for all $n$: for each $n\ge 3$,
the bar cohomology $H^n(\barB_{\mathrm{ch}}(\mathrm{Vir}_c))
\ne 0$ provides nontrivial tree contributions to $\mu_n^M$.

A crucial distinction: \emph{specific} modules may have
$\mu_n^M = 0$ for $n\ge 3$ (for example, the trivial module
$\Bbbk$ with $h'=0$ has $\mu_n^M = 0$ for $n\ge 3$ because
all positive-mode insertions annihilate the vacuum).  The
claim is that for the \emph{category} of all modules, the
$A_\infty$-structure is nontrivial at every arity: there exist
modules $M$ for which $\mu_n^M\ne 0$, for each $n\ge 3$.  The
generic Verma module $M(h',26-c)$ with $h'$ away from the
singular-vector locus provides such modules: the HPL tree
sums for $\mu_n^M$ on generic Verma modules have nonzero
coefficients because the bar cohomology classes at every
arity pair nontrivially with the generic highest-weight
structure.

The physical interpretation: the infinite
$A_\infty$-hierarchy on Virasoro modules encodes the
infinite sequence of null-vector decoupling equations
(BPZ equations) at all levels.  Each $\mu_n^M$ corresponds
to a level-$n$ BPZ differential operator acting on
$n$-point conformal blocks.
\end{proof}

\begin{computation}[Virasoro Ext groups;
\ClaimStatusProvedHere]
% label removed: comp:virasoro-ext
\index{Virasoro algebra!Ext groups}
The Ext groups between Verma modules are:
\begin{equation}
% label removed: eq:virasoro-ext
\operatorname{Ext}^n_{\mathrm{Vir}_{26-c}}
(M(h_1),M(h_2))
\;\cong\;
\begin{cases}
\Bbbk & \text{if } n=0,\;h_1=h_2, \\
0 & \text{if } n=0,\;h_1\ne h_2, \\
\Bbbk^{p(n)} & \text{if } n\ge 1,\;h_1=h_2
  \text{ (generic $c$)},
\end{cases}
\end{equation}
where $p(n)$ is the number of partitions of $n$
(reflecting the graded dimension of the Virasoro algebra
at level~$n$).  At generic central charge $26-c$
(away from the minimal model values): the module category
has non-trivial Ext groups in all degrees, but no
non-trivial extensions between distinct weights.

At minimal model values $c=c_{p,q}=1-6(p-q)^2/(pq)$:
the dual central charge $26-c_{p,q}$ is generically
irrational, and the Ext computation is modified by the
appearance of singular vectors in Verma modules.  The
resulting non-semisimple category is the logarithmic
extension of the $(p,q)$ minimal model.
\end{computation}

\subsection{The tensor product via the $R$-matrix}
% label removed: subsec:tensor-product-r-matrix

The monoidal structure on $\cA^!\text{-}\mathbf{mod}$
is controlled by the spectral $R$-matrix
$r(z)\in\cA^!\hat\otimes\cA^!(\!(z^{-1})\!)$
(Definition~\ref{def:universal-mc-kernel}).

\begin{definition}[Spectral tensor product]
% label removed: def:spectral-tensor-product
\index{tensor product!spectral|textbf}
For $A_\infty$-modules $M$ and $N$ over $\cA^!$, the
\emph{spectral tensor product}
$M\otimes_z N$ is the graded vector space $M\otimes N$
equipped with the $A_\infty$-module structure
\begin{equation}
% label removed: eq:spectral-tensor-product
\mu_n^{M\otimes_z N}(a_1,\ldots,a_{n-1},m\otimes n)
\;:=\;
\sum_{I+J+K=n-1}
\mu_{|I|+1}^M(a_I,m)
\otimes
\mu_{|J|+1}^N(r^{(K)}(z)\cdot a_J,n)
+\text{(higher)},
\end{equation}
where $r^{(K)}(z)$ denotes the $K$-th Taylor coefficient
of the $R$-matrix insertion, and the sum is over ordered
partitions $I\sqcup J\sqcup K=\{1,\ldots,n-1\}$.  The
``higher'' terms involve the higher $A_\infty$-structure
maps $m_k$ on $\cA^!$ and ensure that the
$A_\infty$-module relations
\eqref{eq:ainfty-module-relations} are satisfied on the
tensor product.
\end{definition}

\begin{proposition}[Tensor product satisfies the module
relations; \ClaimStatusProvedHere]
% label removed: prop:tensor-product-module-relations
\index{tensor product!A-infinity module@$A_\infty$-module relations}
The spectral tensor product $M\otimes_z N$
of Definition~\textup{\ref{def:spectral-tensor-product}}
is an $A_\infty$-module over $\cA^!$.  The consistency
is equivalent to the Yang--Baxter equation
(Theorem~\textup{\ref{thm:YBE}}).
\end{proposition}

\begin{proof}
The $A_\infty$-module relations for $M\otimes_z N$
decompose into:
\begin{enumerate}[label=(\roman*)]
\item Terms involving only the $M$-structure maps:
  these vanish by the $A_\infty$-module relations for $M$.
\item Terms involving only the $N$-structure maps:
  these vanish by the $A_\infty$-module relations for $N$.
\item Mixed terms involving both $M$ and $N$ structure
  maps, connected by the $R$-matrix: these vanish if and
  only if $r(z)$ satisfies the classical Yang--Baxter
  equation $[r_{12},r_{13}]+[r_{12},r_{23}]+[r_{13},r_{23}]=0$
  (at leading order) and the full quantum Yang--Baxter
  equation $R_{12}R_{13}R_{23}=R_{23}R_{13}R_{12}$
  (at all orders).
\end{enumerate}
The Yang--Baxter equation is proved in
Theorem~\ref{thm:YBE}, so the $A_\infty$-module relations
hold.
\end{proof}

\begin{computation}[Spectral tensor product for Heisenberg;
\ClaimStatusProvedHere]
% label removed: comp:spectral-tensor-heisenberg
\index{Heisenberg algebra!spectral tensor product}
For $\cF_{q_1}\otimes_z\cF_{q_2}$ over $\cH_{-k}$:
the $R$-matrix is the scalar $r(z)=k\,q_1\,q_2/z$
(Computation~\ref{comp:line-op-heisenberg}), so the
tensor product $\cF_{q_1}\otimes_z\cF_{q_2}
\cong\cF_{q_1+q_2}$
carries the standard Fock module structure with charge
$q_1+q_2$.  The spectral parameter contributes only a
scalar phase: the monodromy around $z=0$ is
$e^{2\pi i\,k\,q_1\,q_2}$.  For $k\in\Z$: the
monodromy is trivial, and the tensor product is
independent of the spectral parameter.
\end{computation}

\begin{computation}[Spectral tensor product for affine
$\mathfrak{sl}_2$; \ClaimStatusProvedHere]
% label removed: comp:spectral-tensor-affine
\index{Kac--Moody!spectral tensor product}
For $W_{\lambda_1}\otimes_z W_{\lambda_2}$ over
$\widehat{\mathfrak{sl}}_2{}_{-k-4}$:
the coproduct is
$\Delta_z(x)=x\otimes 1+1\otimes x
+\sum_{n\ge 1}z^{-n}\Delta_n(x)$,
where $\Delta_n$ encodes the $n$-th pole of the OPE.
At leading order: $\mu_2^{M\otimes_z N}(x,w_1\otimes w_2)
=x\cdot w_1\otimes w_2+w_1\otimes x\cdot w_2
+\Omega(w_1\otimes w_2)/z(k{+}2)+O(z^{-2})$.
The fusion rules (Verlinde coefficients) are recovered:
\[
W_{\lambda_1}\otimes_z W_{\lambda_2}
\;\cong\;
\bigoplus_\lambda
N^{\lambda}_{\lambda_1\lambda_2}\,W_\lambda
\quad
\text{(as $\widehat{\fg}_{-k-4}$-modules, for generic $z$)}.
\]
The decomposition is valid away from poles in $z$; at the
poles, the tensor product acquires logarithmic singularities
corresponding to the non-semisimple Jordan blocks in the
Verlinde category.
\end{computation}


%%%%%%%%%%%%%%%%%%%%%%%%%%%%%%%%%%%%%%%%%%%%%%%%%%%%%%%%%%%%%%%%%%%%%%%%%%%%%
% SECTION — THE YANGIAN AS UNIVERSAL LINE-OPERATOR ALGEBRA
%%%%%%%%%%%%%%%%%%%%%%%%%%%%%%%%%%%%%%%%%%%%%%%%%%%%%%%%%%%%%%%%%%%%%%%%%%%%%
\section{The Yangian as universal line-operator algebra}
% label removed: sec:yangian-universal
\index{Yangian!universal line-operator algebra|textbf}

The dg-shifted Yangian $\Ydg(\cA^!)$ of
Theorem~\ref{thm:Koszul_dual_Yangian} is the universal
algebraic structure controlling line-operator fusion.
This section defines $\Ydg(\cA^!)$ intrinsically from the
convolution algebra, constructs the evaluation
homomorphism, and proves the RTT formalism governs all
fusion.

\subsection{The dg-shifted Yangian from the convolution algebra}
% label removed: subsec:yangian-convolution

\begin{definition}[dg-shifted Yangian from convolution]
% label removed: def:yangian-convolution
\index{dg-shifted Yangian!from convolution algebra|textbf}
Let $\cA$ be a chiral Koszul algebra with modular
convolution dg Lie algebra
$\gAmod=\Convstr(\barB_{\mathrm{ch}}(\cA),\cA)$
(Vol~I, Definition~\ref*{V1-def:modular-convolution-dg-lie}).
The \emph{dg-shifted Yangian} $\Ydg(\cA^!)$ is the
enveloping algebra of the genus-$0$ part of the convolution
algebra, localized at the spectral parameter:
\begin{equation}
% label removed: eq:yangian-convolution-def
\Ydg(\cA^!)
\;:=\;
U\bigl({\gAmod}^{g=0}\bigr)[\![z^{-1}]\!][z],
\end{equation}
where ${\gAmod}^{g=0}$ is the genus-$0$ truncation of the
modular convolution algebra (the summand with loop genus
$g=0$ in the genus decomposition
$\gAmod=\bigoplus_{g\ge 0}{\gAmod}^{(g)}$).

The differential on $\Ydg(\cA^!)$ is
$d_Y=d_{\mathrm{conv}}+[r(z),-]$,
where $d_{\mathrm{conv}}$ is the convolution differential
and $r(z)$ is the genus-$0$ spectral kernel
(Definition~\ref{def:universal-mc-kernel}).  The condition
$d_Y^2=0$ is equivalent to $r(z)$ satisfying the CYBE
(Proposition~\ref{prop:ybe-from-d-squared}).
\end{definition}

\begin{remark}[Yangian versus convolution: strict model
of a homotopy object]
% label removed: rem:yangian-strict-model
\index{dg-shifted Yangian!as strict model}
The dg-shifted Yangian $\Ydg(\cA^!)$ is a strict model
of the homotopy-coherent line-operator algebra: the
full $L_\infty$-convolution algebra
$\Convinf(\barB(\cA),\cA)$ carries higher brackets
$\ell_k$ ($k\ge 3$) encoding homotopy corrections, and
$\Ydg(\cA^!)$ is obtained by passing to the universal
enveloping algebra of the strict dg Lie model $\Convstr$.
The MC moduli of $\Convstr$ and $\Convinf$ coincide
(Convention~\ref{conv:vol2-strict-models}), so the
Yangian correctly parametrizes line-operator deformations.
\end{remark}

\subsection{The evaluation homomorphism}
% label removed: subsec:evaluation-homomorphism

\begin{definition}[Evaluation homomorphism]
% label removed: def:evaluation-homomorphism
\index{evaluation homomorphism|textbf}
For each $A_\infty$-module $M$ over $\cA^!$ and each
$z_0\in\C$, the \emph{evaluation homomorphism}
\begin{equation}
% label removed: eq:eval-hom
\mathrm{ev}_{z_0}\colon
\Ydg(\cA^!)
\;\longrightarrow\;
\End(M)
\end{equation}
is defined by specializing the spectral parameter:
$\mathrm{ev}_{z_0}(T(z)):=T(z)\big|_{z=z_0}$
on the generating matrix, and extending multiplicatively.
This is a dg algebra homomorphism: the RTT relations
at $z=z_0$ reduce to the endomorphism algebra relations
on $\End(M)$.
\end{definition}

\begin{proposition}[Evaluation modules form a dense
subcategory; \ClaimStatusProvedHere]
% label removed: prop:evaluation-modules-dense
\index{evaluation modules!density}
Let $M$ be an $A_\infty$-module over $\cA^!$.  The
\emph{evaluation module} $\mathrm{ev}_{z_0}^*(M)$ (the
pullback of $M$ along $\mathrm{ev}_{z_0}$) is a
$\Ydg(\cA^!)$-module.  The evaluation modules
$\{\mathrm{ev}_{z_0}^*(M)\}_{z_0\in\C,\;M\in\cA^!\text{-mod}}$
generate the category of $\Ydg(\cA^!)$-modules under
tensor products and colimits.
\end{proposition}

\begin{proof}
Every $\Ydg(\cA^!)$-module $N$ is determined by its
spectral decomposition: $N=\bigoplus_{z_0}N_{z_0}$, where
$N_{z_0}$ is the generalized eigenspace of the spectral
parameter.  On each $N_{z_0}$, the Yangian acts through
$\mathrm{ev}_{z_0}$, so $N_{z_0}$ is an evaluation module.
The tensor product of evaluation modules at different
spectral parameters yields indecomposable modules whose
spectral support is the set $\{z_1,\ldots,z_n\}$ of
evaluation points.  By taking colimits (= direct sums
and cokernels), every $\Ydg(\cA^!)$-module is built from
evaluation modules.

The density is strict in the dg sense: the inclusion
functor
$\mathrm{ev}^*\colon
\coprod_{z_0\in\C}\cA^!\text{-}\mathbf{mod}
\hookrightarrow
\Ydg(\cA^!)\text{-}\mathbf{mod}$
induces a quasi-equivalence on the subcategory of
modules with finite spectral support.
\end{proof}

\subsection{The RTT formalism}
% label removed: subsec:rtt-formalism

\begin{theorem}[$\Ydg(\cA^!)$ controls all line-operator
fusion; \ClaimStatusProvedHere]
% label removed: thm:yangian-controls-fusion
\index{RTT relation!controls fusion|textbf}
Let\/ $\cA$ be a logarithmic\/ $\SCchtop$-algebra.  Every fusion process in the
line-operator category $\cC_{\mathrm{line}}$ is computed
by the RTT relation in $\Ydg(\cA^!)$: for any three line
operators $\ell_1,\ell_2,\ell_3$ and any morphism
$f\colon\ell_1\otimes\ell_2\to\ell_3$, there exists
a unique intertwining element
$t_f\in\Ydg(\cA^!)$ such that
\begin{equation}
% label removed: eq:rtt-intertwining
f(v_1\otimes v_2)
\;=\;
\mathrm{ev}_{z_3}(t_f)\cdot
\mathrm{ev}_{z_1}(T_1)\cdot v_1\otimes
\mathrm{ev}_{z_2}(T_2)\cdot v_2,
\end{equation}
and the RTT relation
$R(z_1{-}z_2)\,T_1(z_1)\,T_2(z_2)
=T_2(z_2)\,T_1(z_1)\,R(z_1{-}z_2)$
encodes the consistency of all such fusion morphisms.
\end{theorem}

\begin{proof}
\textbf{Step 1: RTT from bar complex.}
The RTT relation is the statement
$R(u{-}v)\,T_1(u)\,T_2(v)=T_2(v)\,T_1(u)\,R(u{-}v)$
in $\End(V)\otimes\End(V)\otimes\Ydg(\cA^!)[\![u^{-1},v^{-1}]\!]$,
where $V$ is the defining representation and $T(u)$ is
the generating matrix.  This was derived from $d^2=0$ on
the three-fold ordered bar complex in
Construction~\ref{constr:dg-shifted-yangian-from-bar}.

\textbf{Step 2: Fusion as RTT intertwining.}
A fusion morphism $f\colon\ell_1\otimes\ell_2\to\ell_3$
is an intertwiner of $\cA^!$-modules:
$f\in\Hom_{\cA^!}(V_{\ell_1}\otimes V_{\ell_2},V_{\ell_3})$.
Under the identification
$\cA^!\text{-mod}\simeq\Ydg(\cA^!)\text{-mod}$ via
evaluation
(Proposition~\ref{prop:evaluation-modules-dense}), this
intertwiner lifts to a Yangian intertwiner.  The RTT
relation ensures the lift is unique: if
$R(z_1{-}z_2)\,T_1(z_1)\,T_2(z_2)
=T_2(z_2)\,T_1(z_1)\,R(z_1{-}z_2)$,
then composing with $f$ on the right gives
\begin{align*}
f\bigl(R(z)\cdot(v_1\otimes v_2)\bigr)
&=
f\bigl(v_2\otimes v_1\bigr)\text{ (braided)},
\end{align*}
which determines $f$ in terms of the $R$-matrix and
the individual line-operator actions.

\textbf{Step 3: Completeness.}
Every intertwiner of $\Ydg(\cA^!)$-modules between
evaluation modules at $z_1$, $z_2$, $z_3$ is determined
by the RTT algebra: the intertwining space
$\Hom_{\Ydg}(\mathrm{ev}_{z_1}^*M_1\otimes\mathrm{ev}_{z_2}^*M_2,
\mathrm{ev}_{z_3}^*M_3)$
is computed by the centralizer algebra of $T_1(z_1)$
and $T_2(z_2)$ inside $\End(M_3)$, which is the image of
the RTT algebra.  By
Proposition~\ref{prop:evaluation-modules-dense}
(density of evaluation modules), this accounts for all
fusion morphisms.
\end{proof}

\subsection{Virasoro: the gravitational Yangian}
% label removed: subsec:gravitational-yangian

\begin{construction}[Gravitational Yangian;
\ClaimStatusProvedHere]
% label removed: constr:gravitational-yangian
\index{Yangian!gravitational|textbf}
\index{Virasoro algebra!Yangian}
For $\cA=\mathrm{Vir}_c$, the dg-shifted Yangian
$\Ydg(\mathrm{Vir}_{26-c})$ is the
\emph{gravitational Yangian}: the universal algebra
controlling gravitational line-operator fusion.

\emph{Generators.} The generating modes
$T_{(n)}$ ($n\ge 0$) correspond to the Virasoro modes
$L_n$ of the dual algebra $\mathrm{Vir}_{26-c}$:
\begin{equation}
% label removed: eq:gravitational-yangian-generators
T(z)
\;=\;
\sum_{n\ge 0}T_{(n)}\,z^{-n-2}
\;\in\;
\Ydg(\mathrm{Vir}_{26-c})[\![z^{-1}]\!].
\end{equation}
The generating function $T(z)$ is the stress-energy
tensor of the dual Virasoro algebra, viewed as an
element of the Yangian.

\emph{Relations.} The RTT relation for the gravitational
Yangian takes the form
\begin{equation}
% label removed: eq:gravitational-rtt
r_c(z_1{-}z_2)\,T(z_1)\,T(z_2)
-T(z_2)\,T(z_1)\,r_c(z_1{-}z_2)
\;=\;0,
\end{equation}
with the $r$-matrix
\begin{equation}
% label removed: eq:virasoro-r-matrix-explicit
r_c(z)
\;=\;
\frac{\partial T}{z}
+\frac{2T}{z^2}
+\frac{c}{2z^4}.
\end{equation}
Expanding \eqref{eq:gravitational-rtt} in modes:
the $z_1^{-m-2}z_2^{-n-2}$ coefficient is
\begin{equation}
% label removed: eq:gravitational-rtt-modes
[T_{(m)},T_{(n)}]
\;=\;
(m-n)\,T_{(m+n)}
+\frac{26-c}{12}m(m^2-1)\,\delta_{m+n,0},
\end{equation}
which is the Virasoro algebra at central charge $26-c$.
The fourth-order pole in $r_c(z)$ produces the cubic
anomaly term $m(m^2-1)$: this is the conformal anomaly,
not a deformation artifact.

\emph{Coproduct.} The meromorphic coproduct is
\begin{equation}
% label removed: eq:gravitational-coproduct
\Delta_z(T_{(n)})
\;=\;
T_{(n)}\otimes 1+1\otimes T_{(n)}
+\sum_{k=1}^{n+1}\binom{n+1}{k}
\frac{T_{(n-k)}\otimes T_{(k-1)}}{z^{k}}
+\frac{c}{2}\frac{(n+1)n(n-1)}{z^{n+2}}\cdot 1\otimes 1.
\end{equation}
The terms with $z^{-k}$ encode the $k$-th pole of the
$TT$ OPE; the last term (proportional to $c$) is the
contribution from the central charge.
\end{construction}

\begin{remark}[The gravitational Yangian is not a standard Yangian]
% label removed: rem:gravitational-yangian-not-standard
\index{Yangian!gravitational!non-standard}
The gravitational Yangian $\Ydg(\mathrm{Vir}_{26-c})$
differs from the Yangian $Y_\hbar(\fg)$ of a simple Lie
algebra $\fg$ in two respects:
\begin{enumerate}[label=(\roman*)]
\item The $r$-matrix $r_c(z)$ has a fourth-order pole
  at $z=0$ (equation~\eqref{eq:virasoro-r-matrix-explicit}),
  not a simple pole.  Standard Yangians have
  $r(z)\sim\Omega/z$; the Virasoro $r$-matrix has the
  additional terms $2T/z^2+c/(2z^4)$ from the conformal
  anomaly.
\item The algebra $\mathrm{Vir}_{26-c}$ is not a
  finite-dimensional Lie algebra: it is an infinite-dimensional
  graded Lie algebra with generators $L_n$ ($n\in\Z$).
  The Yangian $\Ydg(\mathrm{Vir}_{26-c})$ is therefore
  an infinite-dimensional deformation of $U(\mathrm{Vir}_{26-c}[t])$,
  not a finite-dimensional quantum group.
\end{enumerate}
These departures from the standard Yangian framework are
the gravitational analogues of the gauge-theoretic
Yangian: the fourth-order pole is the gravitational
anomaly (compare the simple pole $\Omega/z$ from the
gauge anomaly), and the infinite-dimensional algebra
reflects the infinite tower of gravitational descendants
$L_{-n}|0\rangle$ ($n\ge 2$).
\end{remark}

\subsection{Affine $\mathfrak{sl}_2$: the quantum group
and $q$-parameter}
% label removed: subsec:affine-quantum-group

\begin{theorem}[$q$-parameter derivation;
\ClaimStatusProvedHere]
% label removed: thm:q-parameter-derivation
\index{quantum group!q-parameter derivation@$q$-parameter derivation|textbf}
For $\cA=\widehat{\mathfrak{sl}}_2{}_k$ with Koszul dual
$\cA^!=Y_\hbar(\mathfrak{sl}_2)$ at
$\hbar=1/(k+2)$, the quantum group
$U_q(\widehat{\mathfrak{sl}}_2)$ arises as the
monodromy representation of the KZ connection, with
$q$-parameter
\begin{equation}
% label removed: eq:q-parameter
q\;=\;e^{i\pi\hbar}\;=\;e^{i\pi/(k+2)}.
\end{equation}
\end{theorem}

\begin{proof}
\textbf{Step 1: KZ connection from the $r$-matrix.}
The KZ connection for $n$ line operators at positions
$z_1,\ldots,z_n\in\C$ is
\begin{equation}
% label removed: eq:kz-connection
\nabla_{\mathrm{KZ}}
\;=\;
d-\hbar\sum_{i<j}
\Omega_{ij}\,d\log(z_i-z_j),
\qquad
\hbar=\frac{1}{k+2},
\end{equation}
where $\Omega_{ij}$ acts as the Casimir
$\Omega=e\otimes f+f\otimes e+\tfrac12 h\otimes h$
on the $i$-th and $j$-th tensor factors.  This is a flat
connection on $\mathrm{Conf}_n(\C)$ valued in
$\End(V^{\otimes n})$ (flatness $\nabla^2=0$ is
the CYBE for $r(z)=\Omega/z$, proved in
Proposition~\ref{prop:ybe-from-d-squared}).

\textbf{Step 2: Monodromy representation.}
The fundamental group
$\pi_1(\mathrm{Conf}_n(\C))$ is the pure braid group
$PB_n$.  The monodromy representation
$\rho_{\mathrm{KZ}}\colon PB_n\to\mathrm{GL}(V^{\otimes n})$
of the KZ connection acts on the elementary braid
$\sigma_i$ (exchanging positions $z_i$ and $z_{i+1}$)
by
\begin{equation}
% label removed: eq:kz-monodromy
\rho_{\mathrm{KZ}}(\sigma_i)
\;=\;
\exp(i\pi\hbar\,\Omega_{i,i+1})
\;=\;
q^{\Omega_{i,i+1}},
\qquad
q=e^{i\pi\hbar}.
\end{equation}
This is the standard Drinfeld--Kohno computation
(\cite{Dri89}, Theorem 3.2; \cite{KL93}): the monodromy
of the logarithmic connection $d-\hbar\,\Omega\,d\log z$
around $z=0$ is $e^{2\pi i\hbar\,\Omega}=q^{2\Omega}$.
The half-monodromy (corresponding to the elementary braid)
is $q^{\Omega}$.

\textbf{Step 3: Quantum group structure.}
The monodromy representation factors through the
quantum group $U_q(\mathfrak{sl}_2)$: the matrix
$R^q=q^{\Omega}$ satisfies the quantum Yang--Baxter
equation
$R^q_{12}R^q_{13}R^q_{23}
=R^q_{23}R^q_{13}R^q_{12}$
(which follows from the CYBE by exponentiation, valid
because $\Omega$ is $\mathfrak{sl}_2$-invariant).  The
quantum group $U_q(\mathfrak{sl}_2)$ is the quotient of
the braid group algebra $\C[B_n]$ by the kernel of
$\rho_{\mathrm{KZ}}$, equipped with the Hopf algebra
structure inherited from the braid group comultiplication.

\textbf{Step 4: Affine extension.}
For the affine algebra
$\widehat{\mathfrak{sl}}_2{}_k$, the loop structure
upgrades $U_q(\mathfrak{sl}_2)$ to
$U_q(\widehat{\mathfrak{sl}}_2)$: the affine root
generators $E_0(u)$, $F_0(u)$ come from the zeroth
root currents of $\widehat{\mathfrak{sl}}_2{}_{-k-4}$,
and the spectral parameter $u$ in the Yangian becomes
the multiplicative spectral parameter in the quantum
affine algebra via $u=e^{i\pi z/(k+2)}$.  The
relation between additive ($z$) and multiplicative ($u$)
spectral parameters is
$u=q^z$, consistent with $q=e^{i\pi/(k+2)}$.
\end{proof}

\begin{computation}[Explicit $q$-parameter for small levels;
\ClaimStatusProvedHere]
% label removed: comp:q-parameter-small-levels
\index{quantum group!q-parameter@$q$-parameter!small levels}
\begin{center}
\renewcommand{\arraystretch}{1.2}
\begin{tabular}{cccl}
\textbf{Level $k$} & $\hbar=1/(k+2)$
  & $q=e^{i\pi/(k+2)}$ & \textbf{Quantum group} \\
\hline
$k=1$ & $1/3$ & $e^{i\pi/3}$
  & $U_{e^{i\pi/3}}(\mathfrak{sl}_2)$:
    root of unity, $q^6=1$ \\[3pt]
$k=2$ & $1/4$ & $e^{i\pi/4}$
  & $U_{e^{i\pi/4}}(\mathfrak{sl}_2)$:
    root of unity, $q^8=1$ \\[3pt]
$k=3$ & $1/5$ & $e^{i\pi/5}$
  & $U_{e^{i\pi/5}}(\mathfrak{sl}_2)$:
    root of unity, $q^{10}=1$ \\[3pt]
$k=-2$ & $\infty$ & $q=1$
  & $U(\mathfrak{sl}_2)$ (classical limit) \\[3pt]
$k\to\infty$ & $0$ & $q\to 1$
  & $U(\mathfrak{sl}_2)$ (large-level limit)
\end{tabular}
\end{center}
At $k=-2$ (critical level): $\hbar=1/(k+2)=\infty$,
the KZ connection degenerates, and the quantum group
reduces to the classical enveloping algebra.  The
Yangian $Y_\hbar(\mathfrak{sl}_2)$ with
$\hbar\to\infty$ is the current algebra
$U(\mathfrak{sl}_2[t])$, consistent with the
Feigin--Frenkel description of the critical-level
affine algebra.
\end{computation}


%%%%%%%%%%%%%%%%%%%%%%%%%%%%%%%%%%%%%%%%%%%%%%%%%%%%%%%%%%%%%%%%%%%%%%%%%%%%%
% SECTION — THE MC3 PROGRAMME AND CATEGORICAL LIFTING
%%%%%%%%%%%%%%%%%%%%%%%%%%%%%%%%%%%%%%%%%%%%%%%%%%%%%%%%%%%%%%%%%%%%%%%%%%%%%
\section{The MC3 programme and categorical lifting}
% label removed: sec:mc3-programme
\index{MC3!programme|textbf}

The MC3 problem is the categorical lift of the
character-level Clebsch--Gordan decomposition to the
full derived module category.  This section states the
problem, records the proved type-$A$ resolution, and
identifies the obstruction to extension.

\subsection{Statement of the MC3 problem}
% label removed: subsec:mc3-statement

\begin{definition}[MC3: categorical Clebsch--Gordan]
% label removed: def:mc3-problem
\index{MC3!problem statement|textbf}
Let $\cA$ be a chiral Koszul algebra with Koszul dual
$\cA^!=Y_\hbar(\fg)$ (the dg-shifted Yangian of a simple
Lie algebra~$\fg$).  The \emph{MC3 problem} is:

\emph{Lift the character-level Clebsch--Gordan
decomposition of $Y_\hbar(\fg)$-modules to a categorical
equivalence in the derived module category
$\Ydg(\fg)\text{-}\mathbf{mod}$.}

Concretely: given two $Y_\hbar(\fg)$-modules $M$ and $N$,
the tensor product $M\otimes_z N$ decomposes into
irreducibles at the character level
(by the Clebsch--Gordan formula for $U_q(\fg)$).
MC3 asks whether this decomposition lifts to a
splitting in the dg module category: is the
tensor product $M\otimes_z N$ quasi-isomorphic to a
direct sum of evaluation modules?
\end{definition}

\begin{remark}[MC3 in the MC hierarchy]
% label removed: rem:mc3-hierarchy
\index{MC frontier!MC3 position}
MC3 is the third of the five Maurer--Cartan levels
(Vol~I, \S\ref*{sec:concordance-mc-frontier}):
\begin{itemize}
\item MC1 (PBW concentration): $H^n(\barB_{\mathrm{ch}}(\cA))
  =0$ for $n\gg 0$.  \textbf{Proved} for all standard
  families.
\item MC2 (bar-intrinsic MC): $\Theta_\cA:=D_\cA-d_0
  \in\mathrm{MC}(\gAmod)$.  \textbf{Proved}
  (Theorem~\ref*{V1-thm:mc2-bar-intrinsic}).
\item MC3 (categorical CG): tensor products of
  $\cA^!$-modules decompose.  \textbf{Proved for all simple types}
  (Vol~I, Corollary~\ref*{V1-cor:mc3-all-types});
  full four-package DK programme proved in type~$A$.
\item MC4 (completion towers): bar-cobar passes to
  inverse limits.  \textbf{Proved}
  (Theorem~\ref*{thm:completed-bar-cobar-strong}).
\item MC5 (all-genera sewing): partition functions converge.
  \textbf{Proved} (Theorem~\ref*{thm:general-hs-sewing}).
\end{itemize}
MC3 is the categorical heart: it converts the spectral
$R$-matrix from a formal object into a genuine
intertwining operator between modules.
\end{remark}

\subsection{Type $A$: the proved resolution}
% label removed: subsec:mc3-type-a

\begin{theorem}[MC3 resolution in type $A$;
\ClaimStatusProvedElsewhere]
% label removed: thm:mc3-type-a-line-ops
\index{MC3!type A resolution|textbf}
\textup{(Vol~I, Theorem~\ref*{thm:mc3-type-a-resolution}.)}
For $\fg=\mathfrak{sl}_N$ ($N\ge 2$) and generic level~$k$
(away from rational values of $q=e^{i\pi/(k+N)}$), the
categorical Clebsch--Gordan decomposition holds:
every finite tensor product of evaluation modules for
$Y_\hbar(\mathfrak{sl}_N)$ decomposes into a direct sum
of evaluation modules in the derived category.
\end{theorem}

This result has been generalized to all simple types
(Vol~I, Corollary~\ref*{V1-cor:mc3-all-types}) via the
multiplicity-free $\ell$-weight approach (Chari--Moura, 2006),
which replaces the minuscule hypothesis used below.

The type-$A$ proof assembles four ingredients.

\subsubsection{The chromatic filtration}
% label removed: subsubsec:chromatic-filtration

\begin{definition}[Chromatic filtration on the Yangian]
% label removed: def:chromatic-filtration
\index{chromatic filtration|textbf}
The \emph{chromatic filtration} on $Y_\hbar(\mathfrak{sl}_N)$
is the weight-space decomposition under the Cartan
subalgebra $\mathfrak{h}\subset\mathfrak{sl}_N$:
\begin{equation}
% label removed: eq:chromatic-filtration
Y_\hbar(\mathfrak{sl}_N)
\;=\;
\bigoplus_{\alpha\in Q}
Y_\hbar(\mathfrak{sl}_N)_\alpha,
\end{equation}
where $Q=\Z\alpha_1\oplus\cdots\oplus\Z\alpha_{N-1}$ is the
root lattice and $Y_\hbar(\mathfrak{sl}_N)_\alpha$ consists
of elements of weight~$\alpha$ (i.e., elements $x$ such
that $[h_i,x]=\alpha(h_i)\cdot x$ for all Cartan generators
$h_i$).

The chromatic filtration refines the standard PBW
filtration: within each weight space, the Yangian generators
$E_i^{(r)}$, $F_i^{(r)}$, $H_i^{(r)}$ have PBW degree~$1$,
and the weight-space decomposition is compatible with the
PBW product.  The associated graded of the chromatic
filtration is the polynomial algebra
$\Bbbk[E_i^{(r)},F_i^{(r)},H_i^{(r)}]_{i,r}$
(by the PBW theorem for Yangians, \cite{CP94}).
\end{definition}

\begin{proposition}[Chromatic filtration compatibility;
\ClaimStatusProvedHere]
% label removed: prop:chromatic-compatibility
\index{chromatic filtration!compatibility}
The chromatic filtration is compatible with the
Yangian coproduct:
$\Delta_z(Y_\alpha)\subset
\bigoplus_{\beta+\gamma=\alpha}Y_\beta\otimes Y_\gamma
[\![z^{-1}]\!][z]$.
In particular, tensor products of weight-space modules
decompose weight-by-weight.
\end{proposition}

\begin{proof}
The coproduct $\Delta_z$ is a homomorphism of
$\mathfrak{h}$-modules (the Cartan subalgebra acts by
derivations on the Yangian, and $\Delta_z$ preserves
this action).  Therefore
$\Delta_z(x)$ has weight $\alpha$ whenever $x$ has
weight $\alpha$, and the tensor product decomposition
$Y_\alpha\otimes Y_\beta\subset Y_{\alpha+\beta}$
is preserved.
\end{proof}

\subsubsection{The prefundamental Clebsch--Gordan closure}
% label removed: subsubsec:prefundamental-cg

\begin{proposition}[Prefundamental Clebsch--Gordan;
\ClaimStatusProvedElsewhere]
% label removed: prop:prefundamental-cg-line-ops
\index{prefundamental representation!Clebsch--Gordan|textbf}
\textup{(Vol~I, Proposition~\ref*{prop:prefundamental-clebsch-gordan}.)}
For $\fg=\mathfrak{sl}_2$, let $V_n$ denote the
$(n+1)$-dimensional evaluation module and $L^-$ the
prefundamental representation of
$Y_\hbar(\mathfrak{sl}_2)$.  Then for all $n\ge 0$:
\begin{equation}
% label removed: eq:prefundamental-cg
V_n\otimes L^-
\;=\;
\bigoplus_{j=0}^{n}L^-(n-2j)
\qquad
\text{(at the character level)},
\end{equation}
where $L^-(s)$ denotes $L^-$ shifted by spectral
parameter~$s$.
\end{proposition}

\begin{proof}[Proof sketch]
The character identity~\eqref{eq:prefundamental-cg} is a
formal power series identity: the character of $V_n\otimes L^-$
is $\chi_{V_n}\cdot\chi_{L^-}=\sum_{j=0}^n\chi_{L^-(n-2j)}$.
This is verified by comparing the $q$-character formulas
on both sides (Vol~I, Proposition~\ref*{prop:prefundamental-clebsch-gordan}; the computation is verified numerically
for $n=1,\ldots,20$ at depth $60$ in the formal parameter).
\end{proof}

\begin{corollary}[Universal character containment;
\ClaimStatusProvedElsewhere]
% label removed: cor:universal-character-containment-line-ops
\index{Verma module!character containment}
\textup{(Vol~I, Corollary~\ref*{cor:universal-character-containment}.)}
For all $\lambda\in\Z_{\ge 0}$:
$\mathrm{ch}(M(\lambda))\le\mathrm{ch}(V_\lambda\otimes L^-)$
(the character of the Verma module is bounded by the character
of the tensor product with the prefundamental).
\end{corollary}

\begin{corollary}[$K_0$-generation; \ClaimStatusProvedElsewhere]
% label removed: cor:k0-generation-line-ops
\index{K-theory@$K_0$!generation}
\textup{(Vol~I, Corollary~\ref*{cor:k0-generation-OY}.)}
Every class $[M(\lambda)]\in K_0(\cO^{\mathrm{sh}})$ lies
in the $K_0$-span of $\{[V_n]\cdot[L^-]\}_{n\ge 0}$.
The $K_0$-level thick generation is complete.
\end{corollary}

\subsubsection{The categorical lift obstruction}
% label removed: subsubsec:categorical-lift-obstruction

\begin{definition}[Categorical CG projectors]
% label removed: def:categorical-cg-projectors
\index{Clebsch--Gordan!categorical projectors|textbf}
The \emph{categorical CG projectors} are the idempotents
$\pi_j\in\End_{Y_\hbar}(V_n\otimes L^-)$ ($0\le j\le n$)
projecting onto the direct summand $L^-(n-2j)$ in the
decomposition~\eqref{eq:prefundamental-cg}.  The projectors
exist at the character level (by the character identity);
the MC3 problem is whether they lift to honest idempotents
in the derived endomorphism algebra
$\operatorname{REnd}_{Y_\hbar}(V_n\otimes L^-)$.
\end{definition}

\begin{proposition}[Categorical lift in type $A$;
\ClaimStatusProvedElsewhere]
% label removed: prop:categorical-lift-type-a
\index{MC3!categorical lift!type A}
\textup{(Vol~I, Theorem~\ref*{thm:mc3-type-a-resolution}.)}
For $\fg=\mathfrak{sl}_N$ and generic~$k$, the categorical
CG projectors exist: the character-level decomposition
lifts to a splitting in the derived category.  The proof
uses:
\begin{enumerate}[label=(\roman*)]
\item \emph{Chromatic filtration}
  (Definition~\textup{\ref{def:chromatic-filtration}}):
  the weight-space decomposition reduces the problem to
  finitely many weight components.
\item \emph{Prefundamental CG closure}
  (Proposition~\textup{\ref{prop:prefundamental-cg-line-ops}}):
  the character-level decomposition provides candidates
  for the categorical projectors.
\item \emph{Efimov completion} (\cite{Efi20}):
  the ind-completion $\mathrm{Ind}\text{-}\cO^{\mathrm{sh}}$
  admits a $t$-structure whose heart contains all the
  summands, allowing the use of Barr--Beck--Lurie monadic
  descent.
\item \emph{Derived Koszul duality on compacts}
  (Vol~I, Theorem~\ref*{thm:dk-fd-typeA}): on compact
  (finite-dimensional) modules, the Koszul duality functor
  is an equivalence, so the CG projectors on compact modules
  lift automatically.  By density of compacts
  (Kaplansky density) and Efimov completion, the lift extends
  to all modules.
\end{enumerate}
\end{proposition}

\subsection{General type: the all-types extension}
% label removed: subsec:mc3-general-type

\begin{theorem}[MC3 for all simple types;
\ClaimStatusProvedElsewhere]
% label removed: conj:mc3-arbitrary-type-line-ops
\index{MC3!all simple types|textbf}
\textup{(Vol~I, Theorem~\ref*{V1-thm:categorical-cg-all-types},
Corollary~\ref*{V1-cor:mc3-all-types}.)}
For any simple Lie algebra $\fg$ and generic level~$k$:
\begin{enumerate}[label=(\roman*)]
\item The multiplicity-free $\ell$-weight property
  (Chari--Moura) gives CG formulas for all Dynkin types,
  replacing the minuscule hypothesis.
\item The categorical CG projectors exist in the derived
  category $D^b(\cO^{\mathrm{sh}}_q(\widehat{\fg}))$.
\item The equivalence
  $\cC_{\mathrm{line}}(\widehat{\fg}_k)
  \simeq\cO^{\mathrm{sh}}_q(\widehat{\fg})$
  is an equivalence of braided monoidal dg categories.
\end{enumerate}
\end{theorem}

\begin{remark}[Resolution via multiplicity-free $\ell$-weights]
% label removed: rem:hernandez-jimbo-obstruction
\index{Hernandez--Jimbo!prefundamental CG}
The original obstruction to extending MC3 beyond type~$A$ was the
prefundamental CG closure for general Dynkin type.  This was
resolved by the multiplicity-free $\ell$-weight approach
(Vol~I, Theorem~\ref*{V1-thm:categorical-cg-all-types}): the
Chari--Moura property~\cite{ChariMoura06} replaces the minuscule
hypothesis with the strictly weaker multiplicity-free $\ell$-weight
property, which holds universally for all simple types.
\end{remark}

\subsection{The DK ladder}
% label removed: subsec:dk-ladder

The Drinfeld--Kohno (DK) programme organizes the relationship
between the Yangian, the quantum group, and the line-operator
category into a sequence of equivalences of increasing
categorical strength.

\begin{definition}[DK levels]
% label removed: def:dk-levels
\index{Drinfeld--Kohno ladder|textbf}
The \emph{DK ladder} is the following sequence of
equivalences, each refining the previous:
\begin{enumerate}[label=\textup{DK-\arabic*.}]
\item \textbf{DK-0} (Laplace shadow):
  $r(z)=\int_0^\infty e^{-\lambda z}\{\cdot_\lambda\cdot\}\,d\lambda$
  identifies the classical $r$-matrix as the Laplace
  transform of the $\lambda$-bracket.
  \textbf{Proved} (Vol~I; Proposition~\ref*{thm:collision-residue-twisting}).
\item \textbf{DK-1} (KZ monodromy = quantum group):
  the monodromy of the KZ connection at level~$k$ is the
  quantum group $U_q(\fg)$ at $q=e^{i\pi/(k+h^\vee)}$.
  \textbf{Proved} (Drinfeld~\cite{Dri89},
  Kazhdan--Lusztig~\cite{KL93}).
\item \textbf{DK-2} (categorical equivalence for integrable
  modules):
  $\cO_k^{\mathrm{int}}(\widehat{\fg})
  \simeq\mathrm{Rep}_q(\fg)$
  (Kazhdan--Lusztig~\cite{KL93}).
  \textbf{Proved} for all types.
\item \textbf{DK-3} (derived equivalence for category
  $\cO$): the derived categories
  $D^b(\cO_k(\widehat{\fg}))
  \simeq D^b(\cO_q(\fg))$
  are equivalent.
  \textbf{Proved} (Arakawa--Bezrukavnikov~\cite{AB09}
  in type $A$ and for generic $q$).
\item \textbf{DK-4} (completed bar-cobar equivalence):
  the completed bar-cobar adjunction gives an equivalence
  on pro-categories.
  \textbf{Proved} by the MC4 strong completion-tower
  theorem (Vol~I, Theorem~\ref*{thm:completed-bar-cobar-strong}).
\item \textbf{DK-5} (full monoidal equivalence):
  the braided monoidal structure on line operators matches
  the braided monoidal structure on quantum group modules,
  including all higher categorical data.
  \textbf{Proved in type $A$} (via MC3 resolution +
  DK-3); MC3 CG closure now proved for all simple types
  (Vol~I, Corollary~\ref*{V1-cor:mc3-all-types}),
  reducing DK-5 for general type to extending DK-3.
\end{enumerate}
\end{definition}

\begin{remark}[DK ladder and the MC hierarchy]
% label removed: rem:dk-mc-relationship
\index{Drinfeld--Kohno ladder!MC relationship}
The DK levels correspond to MC levels:
\begin{center}
\renewcommand{\arraystretch}{1.2}
\begin{tabular}{ll}
\textbf{DK level} & \textbf{MC input} \\
\hline
DK-0 & MC2 (bar-intrinsic MC element) \\
DK-1 & MC2 + analytic continuation \\
DK-2 & MC1 (PBW concentration) + DK-1 \\
DK-3 & MC3 (categorical CG) \\
DK-4 & MC4 (completion towers) \\
DK-5 & MC3 + MC4 + MC5 (all genera)
\end{tabular}
\end{center}
The MC hierarchy provides the algebraic inputs; the DK
ladder translates these into representation-theoretic
equivalences.
\end{remark}


%%%%%%%%%%%%%%%%%%%%%%%%%%%%%%%%%%%%%%%%%%%%%%%%%%%%%%%%%%%%%%%%%%%%%%%%%%%%%
% SECTION — W-ALGEBRA LINE OPERATORS
%%%%%%%%%%%%%%%%%%%%%%%%%%%%%%%%%%%%%%%%%%%%%%%%%%%%%%%%%%%%%%%%%%%%%%%%%%%%%
\section{$\mathcal{W}$-algebra line operators}
% label removed: sec:w-algebra-line-operators
\index{W-algebra@$\mathcal{W}$-algebra!line operators|textbf}

Line operators for $\mathcal{W}_N$-algebras arise from
Drinfeld--Sokolov (DS) reduction of the affine line-operator
category.  The principal case is controlled by the
$\mathcal{W}_N$-Yangian; non-principal reductions produce
a richer structure with obstructions to commutativity
of bar-cobar with DS.

\subsection{Principal $\mathcal{W}_N$: line operators at dual level}
% label removed: subsec:principal-wn-lines

\begin{definition}[$\mathcal{W}_N$ line-operator category]
% label removed: def:wn-line-category
\index{W-algebra@$\mathcal{W}$-algebra!line operator category}
For $\cA=\mathcal{W}_N{}_{c}$ (the principal
$\mathcal{W}$-algebra of type $\mathfrak{sl}_N$ at
central charge~$c$), the line-operator category is
\begin{equation}
% label removed: eq:wn-line-category
\cC_{\mathrm{line}}(\mathcal{W}_N{}_c)
\;\simeq\;
\mathcal{W}_N{}_{c^!_N}\text{-}\mathbf{mod}^{\mathrm{dg}},
\end{equation}
where $c^!_N=2\sum_{s=2}^{N}(6s^2-6s+1)-c$ is the dual
central charge.  The modules are over the Koszul dual
$\cA^!=\mathcal{W}_N{}_{c^!_N}$: the
$\mathcal{W}_N$-algebra at the dual central charge.
\end{definition}

\begin{remark}[Dual central charge]
% label removed: rem:dual-central-charge
\index{W-algebra@$\mathcal{W}$-algebra!dual central charge}
The dual central charge formula
$c^!_N=2\sum_{s=2}^{N}(6s^2-6s+1)-c$ generalizes
the Virasoro duality $c^!_2=26-c$ to higher spin.
For $\mathcal{W}_3$: $c^!_3=100-c$.  For general~$N$:
$c^!_N=c_N^{\mathrm{tot}}-c$, where
$c_N^{\mathrm{tot}}=2\sum_{s=2}^{N}(6s^2-6s+1)$
is the total central charge of the $N-1$ strong generators
of $\mathcal{W}_N$ (a spin-$s$ generator contributes
$2(6s^2-6s+1)$ to the critical central charge).
The self-dual point $c=c_N^{\mathrm{tot}}/2$ is the analogue
of $c=13$ for the Virasoro algebra.
\end{remark}

\subsection{The $\mathcal{W}_3$ case in detail}
% label removed: subsec:w3-detail

\begin{construction}[$\mathcal{W}_3$ Yangian generators
and representation theory; \ClaimStatusProvedHere]
% label removed: constr:w3-yangian
\index{W3@$\mathcal{W}_3$!Yangian generators|textbf}
The $\mathcal{W}_3$-Yangian $\Ydg(\mathcal{W}_3{}_{100-c})$
is generated by two families of modes:
\begin{align}
% label removed: eq:w3-generators
T(z)&=\sum_{n\ge 0}T_{(n)}\,z^{-n-2}
  &&\text{(spin-$2$ current)}, \notag \\
W(z)&=\sum_{n\ge 0}W_{(n)}\,z^{-n-3}
  &&\text{(spin-$3$ current)}.
\end{align}
The spin-$2$ generator $T_{(n)}$ satisfies the Virasoro
relations~\eqref{eq:gravitational-rtt-modes} at central
charge $100-c$; the spin-$3$ generator $W_{(n)}$
satisfies
\begin{equation}
% label removed: eq:w3-tw-commutator
[T_{(m)},W_{(n)}]
\;=\;
(2m-n)\,W_{(m+n)},
\end{equation}
expressing the fact that $W$ is a Virasoro primary of
conformal weight~$3$.  The $WW$ commutation relation is
\begin{equation}
% label removed: eq:w3-ww-commutator
[W_{(m)},W_{(n)}]
\;=\;
(m-n)\Bigl(
\frac{1}{15}(m+n+3)(m+n+2)
-\frac{1}{6}(m+2)(n+2)
\Bigr)\,T_{(m+n)}
+\frac{(100-c)}{360}m(m^2-1)(m^2-4)\,\delta_{m+n,0}
+\beta\,(m-n)\,\Lambda_{(m+n)},
\end{equation}
where $\Lambda_{(n)}=(TT)_{(n)}$ is the normal-ordered
product and $\beta=16/(22+5(100-c))$ is the structure
constant.

\emph{Representation theory.}
The simple objects of
$\mathcal{W}_3{}_{100-c}\text{-}\mathbf{mod}$
are the highest-weight modules $M(h',w')$
parametrized by conformal weight $h'\in\C$ and
$W$-charge $w'\in\C$.  For generic central charge
$100-c$: the Verma modules are irreducible for generic
$(h',w')$, and the module category is semisimple on the
dense open set of parameters.  At special values of
$(h',w')$: singular vectors exist (detected by the
$\mathcal{W}_3$ Kac determinant formula), and the
Verma module has a non-trivial maximal submodule.

The $r$-matrix for $\mathcal{W}_3$ has poles of orders
$1$, $3$, and $5$:
\begin{equation}
% label removed: eq:w3-r-matrix
r_c^{W_3}(z)
\;=\;
\frac{\partial T}{z}
+\frac{2T}{z^2}
+\frac{c}{2z^4}
+\frac{\partial^2 W}{z}
+\frac{5\partial W}{z^2}
+\frac{10 W}{z^3}
+\frac{c_W}{z^6},
\end{equation}
where $c_W$ is the OPE coefficient of the $WW$ central
term.  The fifth-order pole is new compared to the
Virasoro case; it reflects the cubic interaction of the
spin-$3$ current.
\end{construction}

\begin{computation}[$\mathcal{W}_3$ fusion rules at $c=2$;
\ClaimStatusProvedHere]
% label removed: comp:w3-fusion
\index{W3@$\mathcal{W}_3$!fusion rules}
At $c=2$ (dual charge $100-2=98$): the
$\mathcal{W}_3{}_{98}$-module category is semisimple
for generic weights.  The fusion of two modules
$M(h_1,w_1)\otimes_z M(h_2,w_2)$ decomposes into a
direct sum governed by the $\mathcal{W}_3$ Clebsch--Gordan
coefficients, which are computed from the $r$-matrix
\eqref{eq:w3-r-matrix} via the RTT relations.

Under DS reduction from $\widehat{\mathfrak{sl}}_3{}_k$
at $c=c(k)=2$: the three level-$1$
$\widehat{\mathfrak{sl}}_3$-modules ($L_0$,
$L_{\omega_1}$, $L_{\omega_2}$) map to three
$\mathcal{W}_3$-modules with the same fusion ring
$\Z[\Z/3]$ (as in Computation~\ref{comp:ds-on-line-categories}).
\end{computation}

\subsection{Non-principal $\mathcal{W}$-algebras: the hook-type corridor}
% label removed: subsec:non-principal-w-lines

\begin{definition}[Non-principal DS reduction]
% label removed: def:non-principal-ds
\index{Drinfeld--Sokolov reduction!non-principal|textbf}
For a nilpotent element $f\in\mathfrak{sl}_N$ (not
necessarily principal), the Drinfeld--Sokolov reduction
produces the $\mathcal{W}$-algebra
$\mathcal{W}_k(\mathfrak{sl}_N,f)$
(Kac--Roan--Wakimoto~\cite{KRW03}).  The reduction
exists for every nilpotent $f$; the question is whether
bar-cobar duality commutes with the reduction.
\end{definition}

\begin{definition}[Hook-type nilpotent]
% label removed: def:hook-type
\index{hook-type nilpotent|textbf}
\index{W-algebra@$\mathcal{W}$-algebra!hook-type}
A nilpotent $f\in\mathfrak{sl}_N$ is \emph{hook-type}
if its Jordan type is a hook partition
$\lambda=(N-p,1^p)$ for some $0\le p\le N-1$.
Hook-type nilpotents form a $1$-parameter family
interpolating between the zero nilpotent ($p=0$,
$\mathcal{W}=\widehat{\mathfrak{sl}}_N$) and the
principal nilpotent ($p=N-1$, $\mathcal{W}=\mathcal{W}_N$).
\end{definition}

\begin{theorem}[Hook-type Koszul duality;
\ClaimStatusProvedElsewhere]
% label removed: thm:hook-type-koszul
\index{Koszul duality!hook-type|textbf}
\textup{(Fehily \cite{Feh23}; Creutzig--Linshaw--Nakatsuka--Sato
\cite{CLNS24}.)}
For hook-type nilpotents $f=(N-p,1^p)$ in type~$A$:
bar-cobar duality commutes with DS reduction:
\begin{equation}
% label removed: eq:hook-type-koszul
\bigl(\mathcal{W}_k(\mathfrak{sl}_N,f_{(N-p,1^p)})\bigr)^!
\;\simeq\;
\mathcal{W}_{k^\vee}(\mathfrak{sl}_N,f_{(p+1,1^{N-p-1})}),
\end{equation}
where $k^\vee$ is the dual level determined by the duality
involution $f\mapsto f^\vee$ (which sends the hook
$(N-p,1^p)$ to the dual hook $(p+1,1^{N-p-1})$).
\end{theorem}

\begin{construction}[Line-operator category for hook-type
$\mathcal{W}$-algebras; \ClaimStatusProvedHere]
% label removed: constr:hook-type-lines
\index{W-algebra@$\mathcal{W}$-algebra!hook-type!line operators}
The line-operator category for a hook-type
$\mathcal{W}$-algebra is
\begin{equation}
% label removed: eq:hook-type-line-category
\cC_{\mathrm{line}}
\bigl(\mathcal{W}_k(\mathfrak{sl}_N,f_{(N-p,1^p)})\bigr)
\;\simeq\;
\mathcal{W}_{k^\vee}(\mathfrak{sl}_N,f_{(p+1,1^{N-p-1})})
\text{-}\mathbf{mod}^{\mathrm{dg}}.
\end{equation}
This follows from Theorem~\ref{thm:hook-type-koszul}
combined with Theorem~\ref{thm:lines_as_modules}.
The DS reduction functor on line-operator categories
is
\begin{equation}
% label removed: eq:ds-line-functor
\mathrm{DS}_{f}^*\colon
\cO_q^{\mathrm{sh}}(\widehat{\mathfrak{sl}}_N)
\;\longrightarrow\;
\cC_{\mathrm{line}}
(\mathcal{W}_k(\mathfrak{sl}_N,f)),
\end{equation}
sending an $\widehat{\mathfrak{sl}}_N{}_{-k-2N}$-module
$M$ to $H^*_{\mathrm{BRST}}(M\otimes\mathrm{ghost}_f)$,
the BRST cohomology with respect to the nilpotent~$f$.
\end{construction}

\begin{proposition}[Transport propagation for hook-type;
\ClaimStatusProvedHere]
% label removed: prop:transport-propagation
\index{transport propagation!hook-type}
The hook-type Koszul duality
(Theorem~\textup{\ref{thm:hook-type-koszul}}) propagates under
the transport operation: if $(N-p,1^p)$ and $(N-p',1^{p'})$
are two hook-type nilpotents in $\mathfrak{sl}_N$ connected
by an edge in the closure ordering on nilpotent orbits, then
the Koszul duality for $(N-p,1^p)$ implies the Koszul
duality for $(N-p',1^{p'})$ by a spectral sequence argument.
Specifically: the transport map is the composition of
DS reduction from one hook type to the other, and the bar-cobar
commutation at the first hook propagates through the spectral
sequence to the second.
\end{proposition}

\begin{proof}
Let $f\le f'$ in the closure ordering (i.e., the orbit of $f$
is in the closure of the orbit of $f'$).  The DS reduction
$\mathrm{DS}_{f'\to f}\colon
\mathcal{W}_k(\mathfrak{sl}_N,f')\to
\mathcal{W}_{k'}(\mathfrak{sl}_N,f)$
is a functor between vertex algebra categories (it reduces
a larger nilpotent to a smaller one by further BRST
reduction).  On bar complexes, this induces a map
$\barB(\mathrm{DS})\colon
\barB(\mathcal{W}_k(\mathfrak{sl}_N,f'))\to
\barB(\mathcal{W}_{k'}(\mathfrak{sl}_N,f))$.
If the Koszul duality holds for $f'$ (i.e., bar-cobar
commutes with DS at $f'$), then the bar complex map
$\barB(\mathrm{DS})$ is compatible with the Koszul
duality at $f'$, and the induced map on Koszul duals
$(\mathcal{W}_k(\mathfrak{sl}_N,f'))^!
\to(\mathcal{W}_{k'}(\mathfrak{sl}_N,f))^!$
is the DS reduction of the dual algebra.  This is
the ``transport'' of Koszul duality along the closure
ordering.

For hook-type: the hooks $(N-p,1^p)$ ($0\le p\le N-1$)
are totally ordered by the closure relation, and the
transport propagation gives the full hook-type corridor
from the principal case ($p=N-1$, which is the
$\mathcal{W}_N$ duality proved in Vol~I).
\end{proof}

\begin{construction}[Transport-closure in type~$A$]
% label removed: constr:vol2-transport-closure
The transport-closure of hook vertices in the type~$A$ reduction graph $\Gamma_N$ equals the full partition set $\operatorname{Par}(N)$ for all $N \leq 7$ (Vol~I, Theorem~\ref{thm:transport-closure-type-a}).  Every non-hook partition is reachable from a hook vertex via the dominance-order Hasse diagram; the DS--Koszul intertwining propagates along these edges by the transport propagation lemma (Vol~I, Proposition~\ref{prop:transport-propagation}).
\end{construction}

\begin{conjecture}[Type-$A$ transport-to-transpose;
\ClaimStatusConjectured]
% label removed: conj:transport-to-transpose
\index{transport propagation!type A transpose conjecture}
For arbitrary nilpotent $f_\lambda$ in $\mathfrak{sl}_N$
(associated to a partition $\lambda$ of $N$):
\begin{equation}
% label removed: eq:transport-to-transpose
\bigl(\mathcal{W}_k(\mathfrak{sl}_N,f_\lambda)\bigr)^!
\;\simeq\;
\mathcal{W}_{k^\vee_\lambda}(\mathfrak{sl}_N,f_{\lambda^t}),
\end{equation}
where $\lambda^t$ is the transpose partition and
$k^\vee_\lambda$ is the dual level.  This generalizes
the hook-type result~\eqref{eq:hook-type-koszul} to
arbitrary partitions.
\end{conjecture}

\subsection{DS reduction as functor on line-operator categories}
% label removed: subsec:ds-functor-lines

\begin{theorem}[DS as a dg functor on line categories;
\ClaimStatusProvedHere]
% label removed: thm:ds-dg-functor
\index{Drinfeld--Sokolov reduction!dg functor on line categories|textbf}
For the principal nilpotent $f_{\mathrm{prin}}$ in
$\mathfrak{sl}_N$, the DS reduction
$\mathrm{DS}\colon
\widehat{\mathfrak{sl}}_N{}_k\to\mathcal{W}_N{}_{c(k)}$
induces a dg functor on line-operator categories:
\begin{equation}
% label removed: eq:ds-dg-functor
\mathrm{DS}_*\colon
\cO_q^{\mathrm{sh}}(\widehat{\mathfrak{sl}}_N)
\;\longrightarrow\;
\mathcal{W}_N{}_{c^!_N}\text{-}\mathbf{mod}^{\mathrm{dg}},
\end{equation}
compatible with the braided monoidal structure: the functor
preserves tensor products, $R$-matrices, and fusion rules
(up to quasi-isomorphism).
\end{theorem}

\begin{proof}
\textbf{Step 1: DS at the algebra level.}
The DS reduction sends $\widehat{\mathfrak{sl}}_N{}_k$
to $\mathcal{W}_N{}_{c(k)}$ by taking BRST cohomology
with respect to the nilradical $\mathfrak{n}_+$ of a
Borel subalgebra.  On Koszul duals, this induces
$Y_\hbar(\mathfrak{sl}_N)\to\Ydg(\mathcal{W}_N{}_{c^!_N})$
(the DS reduction of the Yangian), which is the
$\mathcal{W}_N$-Yangian.

\textbf{Step 2: DS on modules.}
For a $\widehat{\mathfrak{sl}}_N{}_{-k-2N}$-module $M$:
$\mathrm{DS}_*(M)
=H^*_{\mathrm{BRST}}(M\otimes\mathrm{ghost})$
is a $\mathcal{W}_N{}_{c^!_N}$-module.  The BRST
differential is $Q_{\mathrm{BRST}}=\oint j_{\mathrm{BRST}}(z)\,dz$,
where $j_{\mathrm{BRST}}$ is the BRST current.  Since
$Q_{\mathrm{BRST}}^2=0$ and $Q_{\mathrm{BRST}}$
commutes with the $\mathcal{W}_N$ action (by construction),
$\mathrm{DS}_*$ is a functor of dg categories.

\textbf{Step 3: Monoidal compatibility.}
The tensor product of
$\widehat{\mathfrak{sl}}_N{}_{-k-2N}$-modules is
$M_1\otimes_z M_2$ with the coproduct twisted by
$r(z)=\Omega/z$.  The DS functor sends this to
$H^*_{\mathrm{BRST}}((M_1\otimes_z M_2)\otimes\mathrm{ghost})
\cong
H^*_{\mathrm{BRST}}(M_1\otimes\mathrm{ghost})
\otimes_z
H^*_{\mathrm{BRST}}(M_2\otimes\mathrm{ghost})$
(the K\"unneth isomorphism for BRST cohomology, valid
because the ghost system is a tensor factor).  The
$R$-matrix is preserved because
$Q_{\mathrm{BRST}}$ commutes with the braiding
$R(z)$ (the BRST differential acts only on the
ghost sector, which is decoupled from the spectral
braiding).
\end{proof}

\subsection{The non-principal obstruction}
% label removed: subsec:non-principal-obstruction

\begin{proposition}[Non-principal bar-cobar obstruction;
\ClaimStatusProvedHere]
% label removed: prop:non-principal-obstruction
\index{Drinfeld--Sokolov reduction!non-principal obstruction|textbf}
For non-principal nilpotent $f\ne f_{\mathrm{prin}}$ in
$\mathfrak{sl}_N$ outside the hook-type corridor:
bar-cobar duality does not obviously commute with DS
reduction.  Specifically, the diagram
\begin{equation}
% label removed: eq:non-principal-obstruction-diagram
\begin{tikzcd}
\widehat{\mathfrak{sl}}_N{}_k
  \ar[r, "\barB"]
  \ar[d, "\mathrm{DS}_f"']
& \barB(\widehat{\mathfrak{sl}}_N{}_k)
  \ar[d, "\barB(\mathrm{DS}_f)"] \\
\mathcal{W}_k(\mathfrak{sl}_N,f)
  \ar[r, "\barB"']
& \barB(\mathcal{W}_k(\mathfrak{sl}_N,f))
\end{tikzcd}
\end{equation}
does not commute at the chain level.  The obstruction
is a class
\begin{equation}
% label removed: eq:non-principal-obstruction-class
o_f
\;\in\;
H^1\bigl(
\mathrm{Cone}(\barB\circ\mathrm{DS}_f
\to\mathrm{DS}_f\circ\barB)\bigr)
\end{equation}
measuring the failure of the bar functor to commute with
BRST cohomology.
\end{proposition}

\begin{proof}
The bar construction $\barB$ is a functor on dg algebras
(it sends an algebra to its bar coalgebra).  The DS
reduction $\mathrm{DS}_f$ is a functor on vertex algebras
(it sends a vertex algebra to its BRST cohomology).  The
two functors compose in either order, giving
$\barB\circ\mathrm{DS}_f$ and $\mathrm{DS}_f\circ\barB$
(where, in the second composition, $\mathrm{DS}_f$ acts
on the bar complex by taking BRST cohomology of each
tensor power).

For principal $f$: the BRST differential commutes with
the bar differential because the ghost system is a free
field theory (it is a $bc$-system), and the bar
construction of a free field theory is formal.  For
non-principal $f$: the BRST ghost system involves a
non-abelian gauge constraint (the residual symmetry
group is the Levi factor $L_f\subset\mathfrak{sl}_N$,
which is non-trivial for non-principal $f$), and the
bar construction of this non-abelian system is not formal.
The obstruction class $o_f$ lives in the BRST cohomology
of the mapping cone and vanishes precisely when the
non-abelian gauge constraint is compatible with the
bar differential.

For hook-type nilpotents: the Levi factor is
$\mathfrak{gl}_{N-p}\times\mathfrak{gl}_1^p$, which is
abelian modulo the $\mathfrak{gl}_{N-p}$ factor, and the
obstruction vanishes by the results of
Fehily~\cite{Feh23} and
Creutzig--Linshaw--Nakatsuka--Sato~\cite{CLNS24}.
For general nilpotents: the Levi factor contains
non-abelian simple factors (e.g., for the minimal nilpotent
in $\mathfrak{sl}_4$ with Jordan type $(2,2)$, the Levi
is $\mathfrak{gl}_2$), and the obstruction is expected to
be non-trivial without additional input.
\end{proof}

\begin{remark}[The non-principal frontier]
% label removed: rem:non-principal-frontier
\index{W-algebra@$\mathcal{W}$-algebra!non-principal frontier}
The non-principal obstruction
(Proposition~\ref{prop:non-principal-obstruction}) is the
main open problem in $\mathcal{W}$-algebra Koszul duality.
The hook-type corridor
(Theorem~\ref{thm:hook-type-koszul}) provides the first
genuinely non-principal case where bar-cobar commutes with
DS.  Extension beyond hook-type requires:
\begin{enumerate}[label=(\roman*)]
\item \emph{Understanding the obstruction class} $o_f$:
  for which nilpotents does it vanish?
\item \emph{Transport propagation}
  (Proposition~\ref{prop:transport-propagation}):
  can the hook-type result be propagated to neighboring
  nilpotent orbits in the closure ordering?
\item \emph{The transpose conjecture}
  (Conjecture~\ref{conj:transport-to-transpose}):
  is the Koszul dual always the transpose-nilpotent
  $\mathcal{W}$-algebra?
\end{enumerate}
These questions are active research directions
(Vol~I, \S\ref*{sec:non-principal-w-duality}).
\end{remark}

\begin{remark}[Line operators and the platonic package]
% label removed: rem:lines-platonic-package
\index{platonic package!line operators}
The line-operator category $\cC_{\mathrm{line}}(\cA)
\simeq\cA^!\text{-}\mathbf{mod}$ is the module-theoretic
projection of the platonic package $\Pi_X(L)$
(Vol~I, Construction~\ref*{constr:platonic-package}):
the six-fold datum
$(F_X,\barB_X,\Theta_L,L_L,(V^{\mathrm{br}},T^{\mathrm{br}}),R_4^{\mathrm{mod}})$
encodes the factorization algebra, bar coalgebra, MC element,
modular tangent complex, branch-line reduction, and quartic
resonance class.  The module category is the
representation-theoretic shadow: it sees the evaluation
modules $\mathrm{ev}_z^*(M)$
(Proposition~\ref{prop:evaluation-modules-dense}),
the spectral braiding $R(z)$
(Definition~\ref{def:spectral-braiding}), and the
fusion rules (Construction~\ref{constr:fusion-from-bar}),
but not the full modular structure (genus $\ge 1$
corrections, curvature $\kappa$, shadow Postnikov tower).

The genus-$0$ line-operator data is the tip of the
platonic iceberg.  The full platonic package adds:
\begin{itemize}
\item The genus-$1$ modular correction to the $r$-matrix:
  $R^{\mathrm{mod}}(z;\hbar)=r(z)+\hbar^2 r_1(z)+\cdots$
  (Definition~\ref{def:modular-spectral-kernel}).
\item The shadow Postnikov tower
  $\Theta_\cA^{\le r}$ (Vol~I, Theorem~\ref*{V1-thm:mc2-bar-intrinsic}),
  whose finite-order projections control the $A_\infty$-structure
  on $\cA^!$ at each arity.
\item The modular tangent complex
  (Vol~I, Construction~\ref*{const:vol1-modular-tangent-complex}),
  which governs deformations of the line-operator category
  as the central charge varies.
\end{itemize}
\end{remark}

\begin{computation}[Summary: line-operator categories for
$\mathcal{W}$-algebras; \ClaimStatusProvedHere]
% label removed: comp:w-line-summary
\index{W-algebra@$\mathcal{W}$-algebra!line operator summary}
\begin{center}
\renewcommand{\arraystretch}{1.4}
\begin{tabular}{llll}
\textbf{Algebra} & \textbf{Koszul dual}
  & \textbf{$r$-matrix poles}
  & \textbf{Status} \\
\hline
$\mathcal{W}_2=\mathrm{Vir}_c$
  & $\mathrm{Vir}_{26-c}$
  & $z^{-1},z^{-2},z^{-4}$
  & Proved \\[3pt]
$\mathcal{W}_3{}_c$
  & $\mathcal{W}_3{}_{100-c}$
  & $z^{-1},z^{-2},z^{-3},z^{-4},z^{-5},z^{-6}$
  & Proved \\[3pt]
$\mathcal{W}_N{}_c$
  & $\mathcal{W}_N{}_{c^!_N-c}$
  & $z^{-1},\ldots,z^{-2N}$
  & Proved (principal) \\[3pt]
$\mathcal{W}_k(\fg,f_{\mathrm{hook}})$
  & $\mathcal{W}_{k^\vee}(\fg,f_{\mathrm{hook}}^t)$
  & family-dependent
  & Proved (hook-type) \\[3pt]
$\mathcal{W}_k(\fg,f)$ (general)
  & $\mathcal{W}_{k^\vee}(\fg,f^t)$ (conj.)
  & family-dependent
  & Conjectural
\end{tabular}
\end{center}
The pole structure of the $r$-matrix increases with the
number and spin of the strong generators: a spin-$s$
generator contributes a pole of order $2s$ to the
$r$-matrix.  For $\mathcal{W}_N$: the highest-spin
generator has spin~$N$, giving a maximal pole order
of $2N$.
\end{computation}

\subsection{Planted-forest differential and genus filtration}
% label removed: subsec:planted-forest-genus-filtration

\begin{lemma}[Planted-forest differential and genus shifts]
% label removed: lem:thqg-VII-genus-shifts
\index{planted-forest differential!genus filtration}
Let $\cA$ be a chiral Koszul algebra and let
$\Theta_\cA = \sum_{g,r,d}\hbar^g\,\Theta_{g,r;d}$
be the tridegree decomposition of the bar-intrinsic MC element
\textup{(}Vol~I, Definition~\textup{\ref*{def:vol1-rigid-planted-forest-depth-filtration}}\textup{)}.
The three components of the total differential $D = d_0 + d_{\mathrm{sew}} + d_{\mathrm{pf}}$ act on tridegree as follows:
\begin{enumerate}[label=\textup{(\roman*)},leftmargin=1.8em]
\item The free differential $d_0$ preserves all three indices:
$d_0\colon\Theta_{g,r;d}\to\Theta_{g,r;d}$.
\item The sewing differential $d_{\mathrm{sew}}$ raises the loop
genus $g$ by $1$ and decreases the arity $r$ by $2$ (one
self-gluing):
$d_{\mathrm{sew}}\colon\Theta_{g,r;d}\to\Theta_{g+1,r-2;d}$.
\item The planted-forest differential $d_{\mathrm{pf}}$ raises
the \emph{log boundary depth} $d$ by $1$, and preserves
the loop genus $g$:
$d_{\mathrm{pf}}\colon\Theta_{g,r;d}\to\Theta_{g,r;d+1}$.
\end{enumerate}
In particular, the genus filtration
$F^g = \bigoplus_{g'\ge g}\Theta_{g',\bullet;\bullet}$ is
preserved by $d_{\mathrm{pf}}$: planted forests have genus $0$
\textup{(}they are trees of trees\textup{)}, so
$d_{\mathrm{pf}}$ cannot increase the loop genus.  The only
differential that shifts the genus is $d_{\mathrm{sew}}$.
\end{lemma}

\begin{proof}
Part (i): the free differential $d_0$ acts on each generator
independently, preserving all grading data.

Part (ii): the sewing differential $d_{\mathrm{sew}}$
contracts a self-edge, creating one new loop.  This increases
the first Betti number (= loop genus $g$) by $1$ and consumes
two arity slots (the two endpoints of the sewn edge), giving
$r\mapsto r-2$.  The depth $d$ is unchanged because sewing
is a stable-graph operation, not a log-boundary operation.

Part (iii): the planted-forest differential $d_{\mathrm{pf}}$
arises from the codimension-$1$ boundary strata of the log FM
compactification $\FM_n(X|D)$ (Vol~I,
Construction~\ref*{const:vol1-boundary-operators-residue}).
These strata correspond to planted forests: a tree of collision
clusters, where each cluster is itself a rooted tree.  Since
planted forests are contractible (they are trees), they have
genus $0$ and cannot contribute to the loop genus.  The depth
$d$ increases by $1$ because each application of $d_{\mathrm{pf}}$
adds one level to the planted-forest stratification.
\end{proof}

\begin{remark}[Shadow depth and module complexity]
% label removed: rem:shadow-depth-module-complexity
\index{shadow depth!module complexity}
The shadow depth classification
(Vol~I, \S\ref*{sec:concordance-koszulness-programme})
controls the complexity of the $A_\infty$-module structure:
\begin{center}
\renewcommand{\arraystretch}{1.2}
\begin{tabular}{llll}
\textbf{Class} & $r_{\max}$
  & \textbf{$A_\infty$ depth}
  & \textbf{Example} \\
\hline
G (Gaussian) & $2$ & $\mu_n^M=0$ for $n\ge 3$
  & Heisenberg \\
L (Lie/tree) & $3$ & $\mu_n^M=0$ for $n\ge 4$
  & Affine KM \\
C (contact/quartic) & $4$ & $\mu_n^M=0$ for $n\ge 5$
  & $\beta\gamma$ \\
M (mixed) & $\infty$ & all $\mu_n^M\ne 0$
  & Virasoro, $\mathcal{W}_N$
\end{tabular}
\end{center}
Class~G and L algebras have module categories
that are ``close to classical'' (finitely many non-trivial
higher operations).  Class~M algebras (Virasoro,
$\mathcal{W}_N$) have genuinely infinite $A_\infty$-module
hierarchies: no truncation of the $A_\infty$-structure
captures the full module theory.
\end{remark}
